% !TEX program = xelatex
\documentclass[11pt,openany,fontset=fandol]{ctexbook}
\usepackage[paperwidth=176mm,paperheight=250mm,inner=21mm,outer=18mm,top=20mm,bottom=20mm,headheight=14pt,headsep=6mm,footskip=10mm]{geometry}
\usepackage{amsmath,amssymb,amsthm,bm,mathtools,booktabs,array,enumitem,float}
\usepackage{tikz,fancyhdr,hyperref}
\usetikzlibrary{arrows.meta,calc,patterns,angles,quotes,decorations.markings}
\hypersetup{hidelinks}
\setlength{\parindent}{2em}\setlength{\parskip}{3pt}
\setlength{\emergencystretch}{2em}\linespread{1.18}\raggedbottom
\setcounter{tocdepth}{1}
\ctexset{chapter={name={第,讲},number=\chinese{chapter},format=\LARGE\bfseries\centering,beforeskip=16pt,afterskip=22pt},section={format=\Large\bfseries},subsection={format=\large\bfseries}}
\pagestyle{fancy}\fancyhf{}\fancyhead[LE,RO]{\small\thepage}
\fancyhead[LO]{\small 物理竞赛一轮讲义}\fancyhead[RE]{\small 第六讲\quad 相对论基础}
\renewcommand{\headrulewidth}{0pt}
\fancypagestyle{plain}{\fancyhf{}\fancyfoot[C]{\small\thepage}\renewcommand{\headrulewidth}{0pt}}
\theoremstyle{definition}\newtheorem{example}{例}[chapter]\newtheorem{exercise}{习题}[chapter]
\newenvironment{solution}{\par\medskip\Needspace{5\baselineskip}\noindent\textbf{解}\quad\ignorespaces}{\par\medskip}
\newcommand{\dd}{\mathop{}\!\mathrm d}
\tikzset{axis/.style={->,>=Stealth,line width=.45pt},world/.style={line width=.8pt},light/.style={line width=.6pt,postaction={decorate},decoration={markings,mark=at position .55 with {\arrow{Stealth}}}},ray/.style={line width=.6pt,postaction={decorate},decoration={markings,mark=at position .55 with {\arrow{Stealth}}}},guide/.style={dashed,line width=.35pt},dim/.style={<->,>=Stealth,line width=.35pt},every node/.style={font=\small}}
\newcommand{\fig}[3]{\begin{figure}[H]\centering #3\caption{#2}\label{#1}\end{figure}}
\newcommand{\spacetime}{\draw[axis](-.3,0)--(4.5,0)node[right]{$x$};\draw[axis](0,-.2)--(0,4.6)node[above]{$ct$};\node[below left]at(0,0){$O$};}
\newtheorem{supplement}{补充题}[chapter]
\newenvironment{question}{\par\medskip\noindent\begin{minipage}{\textwidth}}{\end{minipage}\par\medskip}
\newtheorem{extension}{拓展题}[chapter]
\tikzset{world/.style={line width=.8pt}}

\newtheorem{drill}{练习}[chapter]
\newenvironment{drillsolution}[1]{\par\medskip\Needspace{6\baselineskip}\phantomsection\noindent\textbf{练习\ref{#1}详解}\label{sol:#1}\par\smallskip}{\par\medskip}
\newenvironment{exercisesolution}[1]{\par\medskip\Needspace{6\baselineskip}\phantomsection\noindent\textbf{习题\ref{#1}详解}\label{sol:#1}\par\smallskip}{\par\medskip}
\tikzset{normal/.style={dash dot,line width=.45pt},virtual/.style={dashed,line width=.45pt},aux/.style={densely dotted,gray,line width=.35pt},axisline/.style={gray,line width=.3pt}}

% Common visual style for all first-round physics lecture projects.
% Each project includes an identical local copy of this file.
\geometry{paperwidth=176mm,paperheight=250mm,inner=21mm,outer=18mm,top=20mm,bottom=20mm,headheight=14pt,headsep=6mm,footskip=10mm}
\hypersetup{hidelinks}
\setlength{\parindent}{2em}\setlength{\parskip}{3pt}\setlength{\emergencystretch}{2em}
\linespread{1.18}\raggedbottom\setcounter{tocdepth}{1}
\widowpenalty=10000\clubpenalty=10000\displaywidowpenalty=10000
\setlength{\intextsep}{8pt plus1pt minus1pt}
\makeatletter
\g@addto@macro\normalsize{\setlength{\abovedisplayskip}{8pt plus2pt minus2pt}\setlength{\belowdisplayskip}{8pt plus2pt minus2pt}\setlength{\abovedisplayshortskip}{0pt plus2pt}\setlength{\belowdisplayshortskip}{6pt plus2pt minus2pt}}
\makeatother
\providecommand{\Needspace}[1]{}
\newdimen\LectureAvailableSpace
\renewcommand{\Needspace}[1]{\par\penalty-100\begingroup\LectureAvailableSpace=\pagegoal\advance\LectureAvailableSpace by-\pagetotal\ifdim\LectureAvailableSpace<#1\relax\ifdim\pagetotal>0pt\newpage\fi\fi\endgroup}
\renewenvironment{solution}{\par\smallskip\Needspace{3\baselineskip}\noindent\textbf{解}\quad\ignorespaces}{\par\medskip}
\ctexset{chapter={name={第,讲},number=\chinese{chapter},format=\LARGE\bfseries\centering,beforeskip=16pt,afterskip=22pt},section={format=\Large\bfseries},subsection={format=\large\bfseries},subsubsection={format=\normalsize\bfseries,beforeskip=10pt,afterskip=4pt}}
\setlist[enumerate]{leftmargin=2em,itemsep=2pt,topsep=3pt}
\pagestyle{fancy}\fancyhf{}\fancyhead[LE,RO]{\small\thepage}\fancyhead[LO]{\small 物理竞赛一轮讲义}
\renewcommand{\headrulewidth}{0pt}
\fancypagestyle{plain}{\fancyhf{}\fancyfoot[C]{\small\thepage}\renewcommand{\headrulewidth}{0pt}}
\renewenvironment{drillsolution}[1]{\par\Needspace{4\baselineskip}\subsubsection*{练习\ref{#1}详解}\phantomsection\label{sol:#1}\nobreak\noindent\textbf{解}\quad\ignorespaces}{\par\medskip}
\ifdefined\exercisesolution
\renewenvironment{exercisesolution}[1]{\par\Needspace{4\baselineskip}\subsubsection*{习题\ref{#1}详解}\phantomsection\label{sol:#1}\nobreak\noindent\textbf{解}\quad\ignorespaces}{\par\medskip}
\fi
\newcommand{\LectureAnswerSection}{\par\section*{练习解答}\addcontentsline{toc}{section}{练习解答}}
\newcommand{\LectureHeader}[2]{\fancyhead[RE]{\small #1\quad #2}}
\newcommand{\LectureCover}[3]{\begin{titlepage}\centering\vspace*{32mm}{\Huge\bfseries #1\par}\vspace{12mm}{\Large 物理竞赛一轮讲义\par}\vspace{18mm}{\LARGE #2\\[5mm]#3\par}\vfill{\large 2026\par}\end{titlepage}}

\LectureHeader{第六讲}{相对论基础}
\begin{document}
\frontmatter
\LectureCover{光学与近代物理}{第六讲}{相对论基础}
\tableofcontents
\mainmatter\setcounter{chapter}{5}\chapter{相对论基础}
\section{伽利略相对性原理与狭义相对论}

\subsection{惯性系与经典变换}

惯性系是自由质点保持静止或匀速直线运动的参考系。伽利略相对性原理指出，力学规律在所有惯性系中具有相同形式。若 $S'$ 相对 $S$ 沿 $+x$ 以速度 $V$ 匀速运动，且两原点在零时刻重合，经典运动学采用
\[
x'=x-Vt,\quad y'=y,\quad z'=z,\quad t'=t,
\qquad u'_x=u_x-V.
\]
这里的绝对时间使速度直接相减。当 $u_x=c$ 时，式中会给出 $c-V$，不能同时满足所有惯性系中真空光速均为 $c$ 的要求。

狭义相对论采用两条基本原理：所有惯性系中的物理定律具有相同形式；所有惯性系测得的真空光速均为 $c$，与光源的运动无关。以下取 $c=3.00\times10^8\,\mathrm{m/s}$。介质中的光速仍须注明介质的静止系。

光速不变改变的是不同惯性系的时空坐标转换。应用之前，须先明确这些坐标由怎样的测量给出。

低速极限下，相对论变换回到经典变换，说明经典关系在日常低速范围内仍是良好近似。判据涉及相关运动速度与 $c$ 的比值；对于光本身，即使光源或装置运动很慢，也不能先把光速按伽利略规则相减。介质内传播还要明确介质静止系，不能把真空光速不变直接改写为所有参考系中的介质光速相同。

\subsection{事件与时空间隔}

事件由确定的地点和时刻描述，例如粒子产生、炸弹爆炸或闪光到达探测器。一个惯性系用静止尺和分布在各处的同步钟记录坐标，而眼睛接收到光的时刻还包含信号传播时间。改变参考系时，必须变换同一事件的位置和时刻，不能把远处“发生”与眼前“收到”当成一个事件。

两事件的间隔定义为
\begin{equation}\label{eq:hrinterval}
\Delta s^2=c^2\Delta t^2-\Delta x^2-\Delta y^2-\Delta z^2.
\end{equation}
洛伦兹变换保持这一量不变。$\Delta s^2>0$ 为类时间隔，可以由亚光速物体连接；$\Delta s^2=0$ 为类光间隔；$\Delta s^2<0$ 为类空间隔，不能用不超光速的信号连接。
\fig{fig:hrcone}{横轴 $\Delta x$ 与纵轴 $c\Delta t$ 等比例。真空光线斜率为 $\pm1$，类时事件位于光锥内部。}{
\begin{tikzpicture}[scale=.85]
\draw[axis](-3.5,0)--(3.7,0)node[right]{$\Delta x$};\draw[axis](0,-.3)--(0,4.2)node[above]{$c\Delta t$};
\fill[black!6](0,0)--(3.5,3.5)--(-3.5,3.5)--cycle;\draw[light](0,0)--(3.5,3.5);\draw[light](0,0)--(-3.5,3.5);
\node at(0,2.6){未来类时区域};\node at(3,.8){类空};\node at(-3,.8){类空};\node[below left]at(0,0){$O$};
\end{tikzpicture}}

若要把两事件间的时间差解释为随身钟的走时，须先确认它们在同一物体的世界线上。对匀速运动，$\Delta\bm r=\bm u\Delta t$，故
\[
c^2\Delta\tau^2=c^2\Delta t^2-|\Delta\bm r|^2,
\qquad \Delta\tau=\Delta t\sqrt{1-u^2/c^2}.
\]
这是该物体随身钟在两事件间的读数差。若物体在两事件间变速，需沿实际世界线积分 $\tau=\int\sqrt{1-u(t)^2/c^2}\,\dd t$，不能只将首末位移代入匀速公式。图\ref{fig:hrcone}中的类空事件则没有这种由同一物体携带的钟。

比较不同地点的时刻，还需约定怎样对钟。惯性系中的雷达对时可把远处反射事件的时间定义为 $(t_1+t_2)/2$，其中 $t_1,t_2$ 是原点钟记录的发出、返回时刻。在图\ref{fig:hr-sync-detail}中，反射器静止于该系的 $x=D$，去、返真空光程相同，因此
\[
t_B=\frac{t_1+t_2}{2},\qquad D=\frac{c(t_2-t_1)}2.
\]
以后所说的事件时刻均指这样建立的同步钟网所赋予的坐标时间。收到远处消息的时刻还要另计光的传播时间。
\fig{fig:hr-sync-detail}{原点钟的两次读数确定远处反射事件 $B$ 的坐标时间；两轴等比例。}{
\begin{tikzpicture}[scale=.8]
\draw[axis](-.4,0)--(4.5,0)node[right]{$x$};\draw[axis](0,-.4)--(0,6.4)node[above]{$ct$};\draw[world](2.4,0)--(2.4,5.9)node[above]{反射器};
\draw[light](0,.6)--(2.4,3)--(0,5.4);\draw[aux](0,3)--(2.4,3);
\fill(0,.6)circle(1.5pt)node[left]{$ct_1$};\fill(0,5.4)circle(1.5pt)node[left]{$ct_2$};\fill(2.4,3)circle(1.5pt)node[right]{$B(D,ct_B)$};\node[left]at(0,3){$ct_B$};\node[below]at(2.4,0){$D$};
\end{tikzpicture}}

\section{时间、长度与坐标变换}

\subsection{时间膨胀}

若一只钟先后经历两事件，在其静止系中两事件同地，这只钟记录的时间差称固有时间 $\Delta\tau$。实验室中该钟以 $V$ 匀速运动时，
\begin{equation}\label{eq:hrdilation}
\Delta t=\gamma\Delta\tau,\qquad
\gamma=(1-\beta^2)^{-1/2},\qquad\beta=V/c.
\end{equation}
图\ref{fig:hrclock}用横向光钟推导此式。钟系镜距为 $b$，静止周期 $\Delta\tau=2b/c$。实验室的每个半程满足
\[
\left(\frac{c\Delta t}{2}\right)^2
=b^2+\left(\frac{V\Delta t}{2}\right)^2,
\]
整理即得式\eqref{eq:hrdilation}。
\fig{fig:hrclock}{左为钟系的竖直往返，右为实验室的斜向往返。$A,B$ 是同一下镜的发射、回收事件。}{
\begin{tikzpicture}
\draw[very thick](-.6,0)--(.6,0) (-.6,2)--(.6,2);\draw[light](-.1,0)--(-.1,2);\draw[light](.1,2)--(.1,0);\node[below]at(0,0){$A,B$};\node[right]at(.7,1){$b$};
\begin{scope}[shift={(4,0)}]\draw[light](0,0)--(1.5,2);\draw[light](1.5,2)--(3,0);\draw[axisline](0,0)--(1.5,0)--(1.5,2);\node[below]at(0,0){$A$};\node[above]at(1.5,2){$M$};\node[below]at(3,0){$B$};\node[left]at(1.5,1){$b$};\draw[dim](0,-.65)--(3,-.65)node[midway,below]{$V\Delta t$};\end{scope}
\end{tikzpicture}}
光钟证明给出时间膨胀的几何依据；前面由时空间隔得到的式子则作代数核验。两种计算都必须选取同一只钟经历的两个事件。

例如图\ref{fig:hrclock}中的发射事件 $A$ 和回收事件 $B$ 发生在同一下镜上，因而钟系读数差是固有时间。若改选 $A$ 和上镜的反射事件 $M$，两事件在钟系中异地，$t_M-t_A=b/c$ 是坐标时间；$AM$ 本身是光的世界线，间隔为零。使用时间膨胀式前，须明确选的是同一接收器先后经历的哪两个事件。

\subsection{长度收缩}

固有长度 $L_0$ 在物体静止系测量。实验室测运动杆长度 $L$，须在同一实验室时刻记录两端位置。

沿杆发出一束光追赶前镜，再反射回后镜。图\ref{fig:hrlongclock}中，实验室去、返时间为 $L/(c-V)$ 和 $L/(c+V)$，总时间为
\[
\Delta t=\frac{2Lc}{c^2-V^2}.
\]
发射、回收发生在同一后镜，故 $\Delta t=\gamma(2L_0/c)$，得到
\begin{equation}\label{eq:hrlength}
\boxed{L=L_0/\gamma=L_0\sqrt{1-\beta^2}.}
\end{equation}
\fig{fig:hrlongclock}{纵向光钟的两镜世界线。真空光线满足 $|\Delta x|=c\Delta t$，去程时间比返程长；两轴等比例。}{
\begin{tikzpicture}[scale=.55]
\draw[axis](-.3,0)--(8,0)node[right]{$x$};\draw[axis](0,-.3)--(0,8.4)node[above]{$ct$};\draw[world](0,0)--(4.8,8);\draw[world](2.4,0)--(7.2,8);
\draw[light](0,0)--(6,6);\draw[light](6,6)--(4.5,7.5);\node[right]at(6,6){$M$};\node[left]at(4.5,7.5){$B$};\node[below left]at(0,0){$A$};\draw[dim](0,-.7)--(2.4,-.7)node[midway,below]{$L$};
\end{tikzpicture}}
垂直运动方向的长度不变。尺缩是同时测长的结果，照片中各点的光到达相机还存在不同时延，不能直接把照片轮廓当作同时测长。

\fig{fig:hr-simultaneous-length}{同一根杆的两条端点世界线。$AB$ 是实验室同时测量，$AC$ 是杆系同时测量；图示 $V=0.6c$，两轴等比例。}{
\begin{tikzpicture}[scale=.85]
\draw[axis](-.4,0)--(6.4,0)node[right]{$x$};\draw[axis](0,-.4)--(0,5.9)node[above]{$ct$};
\draw[world](0,0)--(3.3,5.5)node[above left]{$x'=0$};\draw[world](2.4,0)--(5.7,5.5)node[above right]{$x'=L_0$};
\coordinate(A)at(1.2,2);\coordinate(B)at(3.6,2);\coordinate(C)at(4.95,4.25);
\draw[aux](.3,2)--(4.3,2);\draw[aux](A)--(C);
\fill(A)circle(1.5pt)node[below left]{$A$};\fill(B)circle(1.5pt)node[below right]{$B$};\fill(C)circle(1.5pt)node[right]{$C$};
\node[above]at(2.4,2){$t\text{ 相同}$};\node[above,rotate=30.964]at(3.2,3.2){$t'\text{ 相同}$};\draw[dim](1.2,1.35)--(3.6,1.35)node[midway,below]{$L$};
\end{tikzpicture}}
图\ref{fig:hr-simultaneous-length}还说明，测量条件变了，选取的端点事件也会改变。对实验室同时事件 $A,B$，有 $\Delta t=0$、$\Delta x=L$。用下文的变换，
\[
\Delta x'=\gamma(\Delta x-V\Delta t)=\gamma L.
\]
两端在杆系中各自静止，故即使 $A,B$ 的杆系时刻不同，它们的坐标差仍为 $L_0$。于是再次得到 $L=L_0/\gamma$。杆系真正同时选取的是 $A,C$；图中的 $C$ 与 $B$ 是前端世界线上的不同事件，不能混为一处。

\subsection{洛伦兹变换与速度变换}

标准构型取 $S'$ 沿 $S$ 的 $+x$ 以 $V$ 运动，两原点零时刻重合。时空均匀性使变换为线性，运动原点满足 $x=Vt$，故可写 $x'=k(x-Vt)$。要求 $x=\pm ct$ 两条光线都变为 $x'=\pm ct'$，得 $t'=k(t-Vx/c^2)$。再由两系互易性及 $V=0$ 的连续极限，求得 $k=\gamma$：
\begin{equation}\label{eq:hrlorentz}
\boxed{x'=\gamma(x-Vt),\qquad
t'=\gamma\left(t-\frac{Vx}{c^2}\right),\quad y'=y,\ z'=z.}
\end{equation}
\fig{o6:boost-axes}{标准构型下的两组坐标轴；虚线沿另一根有色轴的方向投影，光线为 $ct=x$}{\begin{tikzpicture}

\draw[axis](-.3,0)--(5,0)node[right]{$x$};
\draw[axis](0,-.3)--(0,5)node[above]{$ct$};
\draw[light,gray](0,0)--(4.5,4.5)node[above right]{光};
\draw[axis,thick,red](0,0)--(4.5,2.7)node[right]{$x'$};
\draw[axis,thick,blue](0,0)--(2.7,4.5)node[above]{$ct'$};
\coordinate(P)at(2.8,3.6);\fill(P)circle(1.5pt);
\draw[aux](P)--(4.1,3.5);\node[right]at(4.2,3.5){$P(x,ct)$};
\draw[aux](P)--(1,.6);\draw[aux](P)--(1.8,3);
\fill[red](1,.6)circle(1.5pt)node[below right]{$x'_P$};
\fill[blue](1.8,3)circle(1.5pt)node[left]{$ct'_P$};
\node[below left]at(0,0){$O$};

\end{tikzpicture}}
图\ref{o6:boost-axes}取 $V=0.6c$，黑色两轴等比例。$ct'$ 轴是飞船原点的世界线 $x=Vt$；$x'$ 轴则由 $t'=0$ 给出 $ct=(V/c)x$。一般的 $t'=\text{常量}$ 线与 $x'$ 轴平行，$x'=\text{常量}$ 线与 $ct'$ 轴平行。因此读出同一事件 $P$ 的有色坐标，要沿另一根有色轴作平行线；不能向黑色轴作垂直投影后直接当作 $x',ct'$。有色轴上的单位刻度还须按洛伦兹变换标定，这不是保持欧氏距离的普通旋转。

这个系数可展开一步求出。互易性要求逆变换为 $x=k(x'+Vt')$，把正变换代入便有
\[
x=k^2\left[x-Vt+V\left(t-\frac{Vx}{c^2}\right)\right]
=k^2(1-\beta^2)x.
\]
所以 $k^2(1-\beta^2)=1$，取与 $V=0$ 连续的正根即为 $\gamma$。将两式代入 $c^2\Delta t'^2-\Delta x'^2$，交叉项相消，剩下
\[
c^2\Delta t'^2-\Delta x'^2
=\gamma^2(1-\beta^2)(c^2\Delta t^2-\Delta x^2)
=c^2\Delta t^2-\Delta x^2.
\]
这既验证了间隔不变，也可用来检查数值变换是否把同一对事件选对。

逆变换把 $V$ 换成 $-V$。坐标差也满足同样关系。若 $\Delta t=0$，通常有 $\Delta t'=-\gamma V\Delta x/c^2\ne0$，可见不同位置的同时性取决于参考系。

对逆变换微分，再除以 $\dd t$，得
\begin{equation}\label{eq:hrvelocity}
u_x=\frac{u'_x+V}{1+Vu'_x/c^2},\qquad
u_y=\frac{u'_y}{\gamma(1+Vu'_x/c^2)},\quad
u_z=\frac{u'_z}{\gamma(1+Vu'_x/c^2)}.
\end{equation}
速度式的分母来自时间变换。沿粒子相邻事件，有 $\dd x'=u'_x\dd t'$，故
\[
\dd t=\gamma\left(1+\frac{Vu'_x}{c^2}\right)\dd t',\qquad
\dd x=\gamma(u'_x+V)\dd t'.
\]
相除得到纵向速度式；再用 $\dd y=\dd y'$、$\dd z=\dd z'$，便得到横向速度式。横向坐标虽不变，横向速度仍会因时间间隔的变化而改变。

把三个速度分量合并，还可得到
\[
1-\frac{u^2}{c^2}
=\frac{(1-V^2/c^2)(1-u'^2/c^2)}{(1+Vu'_x/c^2)^2}.
\]
右端在 $u'<c,V<c$ 时为正，所以合成速度小于 $c$；对光有 $u'=c$，两边同时为零。这里比较的是速度矢量的总大小，不应只检验某一个分量。

若由 $S$ 变到 $S'$，将 $V$ 换为 $-V$。对真空光代入 $u=c$，合成后总光速仍为 $c$；亚光速物体的相对速度仍小于 $c$。同一实验室中两个物体的间距减少率可以大于 $c$，它不是某一物体在另一物体静止系中测得的速度。

\fig{fv:o6:simultaneity}{同时线随参考系倾斜，测量运动长度必须重选事件}{\begin{tikzpicture}[scale=.95]
\draw[->](-.5,0)--(5,0)node[right]{$x$};\draw[->](0,-.3)--(0,3.7)node[above]{$ct$};
\draw[thick](1,-.2)--(1,3.3);\draw[thick](3.8,-.2)--(3.8,3.3);
\draw[blue,thick](.3,1)--(4.5,1)node[right]{$t=\mathrm{const}$};
\draw[red,thick](.3,.65)--(4.5,2.75)node[right]{$t'=\mathrm{const}$};
\fill(1,1)circle(1.5pt)node[above left]{$A$};\fill(3.8,1)circle(1.5pt)node[below right]{$B$};
\fill[red](3.8,2.4)circle(1.5pt)node[above left]{$B'$};
\end{tikzpicture}}
图\ref{fv:o6:simultaneity}把杆画成在 $S$中静止的两条端点世界线。蓝线选出的 $A,B$满足 $\Delta t=0$；变换后却有 $\Delta t'=-\gamma V\Delta x/c^2\ne0$，不能用这对事件直接测量 $S'$中的杆长。$S'$的同时线满足 $ct=(V/c)x+\text{常量}$，应另选同一红线上两端的事件。长度收缩中的“同时测两端”，因而是选事件的条件，不是对同一对任意事件只作空间坐标变换。


图\ref{fv:o6:simultaneity}中杆在 $S$ 静止，故 $x_B-x_A=x_{B'}-x_A=L_0$。红线选取 $A,B'$，满足 $\Delta t'=0$，所以 $\Delta t=VL_0/c^2$。将它代入空间变换：
\[
 L'=\Delta x'=\gamma(L_0-V\Delta t)
 =\gamma L_0(1-\beta^2)=\frac{L_0}{\gamma}.
\]
这与前面“杆在 $S'$ 静止、在 $S$ 运动”的情形互为参照。每一方都在自己的同时线上选端点，所比较的是不同事件对，因而两方都测得对方的运动杆收缩并不矛盾。


\subsection{事件选择与测量条件}

每个测量量都对应特定的事件选择和参考系。时间膨胀选同一只钟的两次读数，长度收缩选测量系中同一时刻的两端，速度变换选同一粒子的相邻事件。对于爆炸、遮光、碰撞这类单个事件，应同时变换它的空间和时间坐标。
\begin{center}\small
\begin{tabular}{p{.20\textwidth}p{.32\textwidth}p{.34\textwidth}}
\toprule
所求量&所选事件或条件&可用关系\\\midrule
运动钟走时&同一随身钟经历两个事件&$\Delta t=\gamma\Delta\tau$\\
运动杆长度&测量系同时记录两端&$L=L_0/\gamma$\\
粒子速度&同一粒子上的相邻事件&$\dd x/\dd t$ 的洛伦兹变换\\
远处事件坐标&位置与时刻属于同一事件&同时变换 $x,t$\\
\bottomrule
\end{tabular}
\end{center}
例\ref{ex:hr5}和例\ref{ex:hr6}的不同“同时面”，以及练习\ref{hw:hr1}中跑道长度与赛跑位移的区别，都可先由这张表确定事件，再进行计算。

\section{两个光学实验}

\subsection{斐索实验与运动介质中的光速}

图\ref{fig:hrfizeau}表示干涉装置中两个等长水管的传播部分。两束光绕行方向相反，一束在两管内均顺水，另一束均逆水。令每管长度为 $l$，水速为 $V\ll c$，水静止系中的光速为 $c/n$，忽略色散。
\fig{fig:hrfizeau}{两管中的水流方向相反。两条互逆光路分别累计 $2l$ 的顺水路程和 $2l$ 的逆水路程。}{
\begin{tikzpicture}[scale=.85]
\fill[black!5](0,.6)rectangle(5,1.6);\fill[black!5](0,-1.6)rectangle(5,-.6);\draw(0,.6)rectangle(5,1.6);\draw(0,-1.6)rectangle(5,-.6);
\draw[ray](0,1)--(5,1);\draw[ray](5,1.3)--(0,1.3);\draw[ray](5,-1)--(0,-1);\draw[ray](0,-1.3)--(5,-1.3);
\draw[axis](1.5,2)--(3.5,2)node[midway,above]{$V$};\draw[axis](3.5,-2)--(1.5,-2)node[midway,below]{$V$};
\draw[aux](5,1)--(5.8,1)--(5.8,-1)--(5,-1);\draw[aux](0,1)--(-.8,1)--(-.8,-1)--(0,-1);
\node[left]at(-.8,0){分束与合束};\draw[dim](0,.15)--(5,.15)node[midway,above]{$l$};
\end{tikzpicture}}
对顺水传播，由前述速度变换得到
\[
u_+=\frac{c/n+V}{1+V/(nc)}
\simeq\frac cn+V\left(1-\frac1{n^2}\right).
\]
逆水速度的大小为 $|u_-|\simeq c/n-V(1-1/n^2)$。若以 $k=1-1/n^2$ 表示一阶系数，两束光的时间差为
\[
\Delta t=\frac{2l}{c/n-kV}-\frac{2l}{c/n+kV}
\simeq\frac{4lV(n^2-1)}{c^2}.
\]
该结果由介质系与装置系之间的速度变换得到，无须假设一种绝对静止的传播介质。

一阶展开中的两项都须保留。由 $(1+V/(nc))^{-1}\simeq1-V/(nc)$，
\[
\left(\frac cn+V\right)\left(1-\frac V{nc}\right)
=\frac cn+V-\frac V{n^2}+O(V^2/c).
\]
可见直接把 $c/n$ 与 $V$ 相加会漏掉 $-V/n^2$。在 $n\to1$ 的真空极限，两方向的速率都回到 $c$；在 $V\to0$ 时，则回到静止介质中的 $c/n$。这两个极限也能检查拖曳系数的形式。

\subsection{迈克耳孙—莫雷实验}

图\ref{fig:hrmichelson}中，两臂在装置静止系长度均为 $l$，分束镜 $M$ 将光分到 $M_1,M_2$，反射后在 $M$ 合束。
\fig{fig:hrmichelson}{等臂干涉仪。右臂平行于所假设的绝对运动方向，上臂与其垂直。}{
\begin{tikzpicture}
\draw[very thick](-.35,-.35)--(.35,.35);\draw[very thick](-.6,3)--(.6,3) (3,-.6)--(3,.6);
\draw[light](-2,0)--(0,0);\draw[light](.08,0)--(.08,3);\draw[light](-.08,3)--(-.08,0);\draw[light](0,.08)--(3,.08);\draw[light](3,-.08)--(0,-.08);\draw[light](0,0)--(0,-1.5);
\node[below left]at(0,0){$M$};\node[above]at(0,3){$M_1$};\node[right]at(3,0){$M_2$};\node[left]at(-2,0){$S$};\node[below]at(0,-1.5){$T$};
\draw[dim](-.9,0)--(-.9,3)node[midway,left]{$l$};\draw[dim](0,-.8)--(3,-.8)node[midway,below]{$l$};\draw[axis](1.2,2)--(3.2,2)node[right]{$V$};
\end{tikzpicture}}
按“相对绝对静止媒质光速为 $c$，装置作伽利略运动”的假设，平行臂与垂直臂的往返时间分别为
\[
t_\parallel=\frac l{c-V}+\frac l{c+V}
=\frac{2l/c}{1-\beta^2},\qquad
t_\perp=\frac{2l/c}{\sqrt{1-\beta^2}},\quad \beta=V/c.
\]
低速展开得 $t_\parallel-t_\perp\simeq l\beta^2/c$。装置转过 $90^\circ$ 后差值反号，预测的条纹移动为
\[
\Delta N\simeq\frac{2l}{\lambda}\frac{V^2}{c^2}.
\]
取 $l=11\,\mathrm m,V=3.0\times10^4\,\mathrm{m/s},\lambda=5.9\times10^{-7}\,\mathrm m$，预测约 $0.37$ 条。实验未出现这一预期的方向效应。狭义相对论在装置静止系直接给出两臂往返时间均为 $2l/c$，理想等臂装置转动后不会产生上述条纹移动。

上述预测是在一组特定假设下计算的：真空光相对某个绝对静止媒质各向同性，而装置的长度与时间又按经典规则处理。实验未出现预期的旋转效应，约束的是这组假设的共同预测，不能从结果单独推出“两束光都没有传播时间”或“仪器不发生干涉”。

图\ref{fig:hrmichelson}中，两臂往返相位差才决定条纹位置。旋转前若有固定的臂长误差，可能产生一个固定相位差；检验所关心的是旋转后随方向改变的那部分。在装置静止系使用相对论时，两臂的光速统一为 $c$，理想等臂的方向依赖自然消失，计算条件和观测量保持一致。

\section{动量、质量与能量}

\subsection{相对论动量与动能}

以下统一用 $m_0$ 表示粒子的不变质量，记其速度为 $u$，$\gamma_u=(1-u^2/c^2)^{-1/2}$。相对论动量、总能量为
\begin{equation}\label{eq:hrenergy}
\bm p=\gamma_um_0\bm u,\qquad E=\gamma_um_0c^2.
\end{equation}
若将 $E/c^2$ 另记为“动质量” $m_{\text{动}}$，则 $m_{\text{动}}=\gamma_um_0$，$\bm p=m_{\text{动}}\bm u$。这只是另一种记号；计算合成粒子的质量时，须明确所求的是不变质量，不能把各粒子的不变质量直接相加。

力定义为 $\bm F=\dd\bm p/\dd t$，功满足 $\dd K=\bm F\cdot\dd\bm r=\bm u\cdot\dd\bm p$。动能的计算以静止时 $K=0$为边界，积分结果记为
\begin{equation}\label{eq:hrkinetic}
\boxed{K=(\gamma_u-1)m_0c^2=E-m_0c^2.}
\end{equation}
对沿直线的运动，$p=m_0u/\sqrt{1-u^2/c^2}$，先对 $\gamma_u$ 使用链式法则：
\[
 \frac{\dd\gamma_u}{\dd u}
 =-\frac12(1-u^2/c^2)^{-3/2}\frac{-2u}{c^2}
 =\frac{u\gamma_u^3}{c^2}.
\]
再对 $p=m_0u\gamma_u$ 使用乘积法则，
\begin{align*}
 \frac{\dd p}{\dd u}
 &=m_0\left(\gamma_u+\frac{u^2}{c^2}\gamma_u^3\right)\\
 &=m_0\gamma_u^3\left(1-\frac{u^2}{c^2}+\frac{u^2}{c^2}\right)
 =m_0\gamma_u^3.
\end{align*}
于是 $\dd K=u\dd p=m_0u\gamma_u^3\dd u=m_0c^2\dd\gamma_u$。
静止时 $u=0$、$\gamma_u=1$、$K=0$，积分上下限一一对应：
\[
 \int_0^{K(u)}\dd K'
 =m_0c^2\int_1^{\gamma_u}\dd\gamma',
 \qquad K(u)=m_0c^2[\gamma']_1^{\gamma_u}
 =(\gamma_u-1)m_0c^2.
\]
这样得到式\eqref{eq:hrkinetic}，积分常量已由静止时的动能零点确定。同一微分关系还给出纵向运动的
\[
F=m_0\gamma_u^3\frac{\dd u}{\dd t}.
\]
因而实验室恒力并不对应恒加速度。粒子越快，同样的动量增量所造成的速度增量越小；继续做功主要使总能量和动量增长。

\fig{o6:momentum-energy-curves}{实线为相对论关系，虚线为经典近似；各图保持粒子不变质量相同}{\begin{tikzpicture}

\begin{scope}[x=4cm,y=.9cm]
\draw[axis](0,0)--(1.08,0)node[right]{$u/c$};
\draw[axis](0,0)--(0,3.4)node[above]{$p/(m_0c)$};
\draw[blue,thick,domain=0:.94,samples=100]plot(\x,{\x/sqrt(1-(\x)^2)});
\draw[dashed](0,0)--(.94,.94);
\node[below]at(1,0){$1$};\draw[aux](1,0)--(1,3.1);
\node[blue,left]at(.9,2.6){相对论};\node[above]at(.58,.58){经典};
\end{scope}
\begin{scope}[xshift=6.8cm,x=4cm,y=1.2cm]
\draw[axis](0,0)--(1.08,0)node[right]{$u/c$};
\draw[axis](0,0)--(0,2.55)node[above]{$K/(m_0c^2)$};
\draw[red,thick,domain=0:.94,samples=100]plot(\x,{1/sqrt(1-(\x)^2)-1});
\draw[dashed,domain=0:.94,samples=80]plot(\x,{.5*(\x)^2});
\node[below]at(1,0){$1$};\draw[aux](1,0)--(1,2.3);
\node[red,left]at(.9,1.7){相对论};\node[above]at(.63,.2){经典};
\end{scope}

\end{tikzpicture}}
图\ref{o6:momentum-energy-curves}中，两种关系在 $u/c\ll1$ 时接近；接近光速时，实线的动量和动能都急剧增大。对有质量粒子，$u\to c^-$ 要求无界的能量，有限做功不能使它达到或越过 $c$。图中低速近似的抛物线不能向高速无限外推。

低速展开 $\gamma_u\simeq1+u^2/(2c^2)$，才退化为 $K\simeq m_0u^2/2$。恒力对应动量随时间线性增长，速度却不会随时间无限增大，例\ref{ex:hr11}将具体计算。

\subsection{换参考系时的能量与动量}

总能量、动量分别依赖参考系，但同一系统的不变质量不依赖参考系。设 $S'$ 相对 $S$ 沿 $+x$ 以速度 $V$ 运动。将速度变换代入 $E=\gamma_um_0c^2$、$p_x=\gamma_um_0u_x$，利用
\[
 \gamma_{u'}=\gamma_V\gamma_u\left(1-\frac{Vu_x}{c^2}\right),
\]
得到
\[
 E'=\gamma_V(E-Vp_x),\qquad
 p_x'=\gamma_V\left(p_x-\frac{VE}{c^2}\right),\qquad
 p_y'=p_y,\quad p_z'=p_z.
\]
这些关系是线性的，对各粒子相加后仍成立，所以也可变换整个碰撞系统的总能量与总动量。不能把 $S$ 系的动能与 $S'$ 系的动量放入同一质量公式，也不能只变换入射粒子而让靶粒子在两系中都保持静止。

若在 $S$ 系有 $E_{\text{初}}=E_{\text{末}}$、$\bm P_{\text{初}}=\bm P_{\text{末}}$，代入上述线性关系便同时得到 $S'$ 系的两条守恒式。守恒不表示数值在不同系中相同，而是同一系内的初、末总量相同。比较一个碰撞前后时，先固定参考系，再逐项列出参与物体；需要换系时，将所有参与物体一起变换。

\subsection{能量动量关系与守恒}

消去 $\gamma_u,u$，得到
\begin{equation}\label{eq:hrem}
\boxed{E^2-p^2c^2=m_0^2c^4,\qquad \bm u=\frac{c^2\bm p}{E}.}
\end{equation}
式\eqref{eq:hrem}可直接从同一粒子的能量、动量求出：
\[
 E^2-p^2c^2=\gamma_u^2m_0^2c^4(1-u^2/c^2)=m_0^2c^4.
\]
\fig{o6:mass-hyperbola}{同一有质量粒子在不同参考系中的能量、动量数值落在同一正能量双曲线上}{\begin{tikzpicture}

\begin{scope}[x=1.1cm,y=1.1cm]
\draw[axis](-2.8,0)--(3,0)node[right]{$p_x/(m_0c)$};
\draw[axis](0,-.2)--(0,3.4)node[above]{$E/(m_0c^2)$};
\draw[gray,dashed](-2.7,2.7)--(0,0)--(2.7,2.7);
\draw[thick,domain=-2.5:2.5,samples=120]plot(\x,{sqrt(1+(\x)^2)});
\fill[blue](1.732051,2)circle(1.6pt)node[right]{$S:(\sqrt3,2)$};
\fill[red](0,1)circle(1.6pt)node[below left]{$S':(0,1)$};
\draw[axis,blue](1.45,1.95)..controls(.9,1.4)and(.55,1.22)..(.2,1.12);
\node[above]at(0,4.2){$E^2-p_x^2c^2=m_0^2c^4$};
\end{scope}

\end{tikzpicture}}
图\ref{o6:mass-hyperbola}只考虑沿 $x$ 的运动。在 $S$ 中取 $E=2m_0c^2,p_x=\sqrt3m_0c$，则粒子速度为 $u=\sqrt3c/2$；选 $V=u$ 变到粒子静止系后，$p'_x=0,E'=m_0c^2$。能量与动量都改变，双曲线所代表的不变质量保持不变。灰虚线对应 $E=|p_x|c$，有质量粒子的数值位于其上方。

光子的 $m_0=0$，故 $E=pc=h\nu=hc/\lambda$。光子有能量和动量，没有可供选取的静止参考系。

孤立系统的总能量、总动量守恒。把整个系统视为一个整体，其不变质量 $M$ 由
\[
M^2c^4=E_{\text{总}}^2-c^2|\bm p_{\text{总}}|^2
\]
确定。完全非弹性碰撞中，部分相对运动能转化为复合物的内能，所以 $M$ 可以大于碰前各粒子不变质量之和。衰变、光子反冲和康普顿散射也应同时使用能量、动量守恒。

守恒式中的系统边界须一致。若末态有复合粒子和辐射，二者都应列入：
\[
\sum E_{\text{初}}=E_{\text{复合}}+E_{\text{辐射}},\qquad
\sum\bm p_{\text{初}}=\bm p_{\text{复合}}+\bm p_{\text{辐射}}.
\]
只有初态总动量为零且末态没有外逸动量时，复合粒子才在该系静止，才能直接以 $Mc^2$ 表示它的全部能量。若总动量 $\bm P\ne0$，则先由 $E_{\text{总}}^2-P^2c^2$ 求系统不变质量；若 $M>0$，系统整体静止系相对当前系的速度为 $\bm V_{\rm cm}=c^2\bm P/E_{\text{总}}$，其大小严格小于 $c$。当 $M=0$ 且总能量非零时，该表达式的大小为 $c$，不对应可选取的惯性静止系。

\fig{fig:hr-photon-system}{两个光子的单独静质量均为零，合系统的静质量仍取决于夹角。左图反向等能，总动量为零；右图同向，总动量为 $2E_\gamma/c$。}{
\begin{tikzpicture}[scale=.9]
\fill(0,0)circle(1.5pt);\draw[light](0,0)--(-2.3,0)node[above]{$E_\gamma$};\draw[light](0,0)--(2.3,0)node[above]{$E_\gamma$};\node[below]at(0,-.55){$M=2E_\gamma/c^2$};
\begin{scope}[shift={(6.6,0)}]\draw[light](0,.3)--(2.8,.3)node[right]{$E_\gamma$};\draw[light](0,-.3)--(2.8,-.3)node[right]{$E_\gamma$};\node[below]at(1.4,-.8){$M=0$};\end{scope}
\end{tikzpicture}}
图\ref{fig:hr-photon-system}可由 $E_\gamma=pc$ 直接验证。一般两个光子有
\[
M^2c^4=2E_1E_2(1-\cos\theta).
\]
同向时整体仍没有静止系；反向等能时则有静止系，能量全部计入系统静能。这说明“各粒子的静质量相加”与“求整个系统的静质量”是不同的运算。练习\ref{hw:hr5}将用这一关系反推母粒子的质量。

\section{非惯性系、局部引力与时钟}

\subsection{惯性力与等效原理}

在低速近似下，相对于惯性系做平动加速的参考系中，可以引入惯性力 $\bm F_{\mathrm{in}}=-m_{\mathrm i}\bm a$，使运动方程保持通常的形式。引力中的 $m_{\mathrm g}$ 决定物体受到的引力，惯性质量 $m_{\mathrm i}$ 则决定物体对加速的响应。自由落体加速度与物体组成无关，表明二者成比例；采用通常的单位约定可令 $m_{\mathrm g}=m_{\mathrm i}$。

图\ref{fig:hr-gr1}比较一个向上加速的小舱和一个处于引力场中的小舱。对足够小的空间区域和足够短的时间，舱内实验不能区分均匀引力与相应的加速效应。这是等效原理的局部表述。
\fig{fig:hr-gr1}{小舱内的加速效应与局部引力效应}{
\begin{tikzpicture}[scale=.9]
\draw[world](0,0)rectangle(3,3);\fill(1.5,1.5)circle(3pt);\draw[axis](1.5,1.4)--(1.5,.5)node[right]{$m\bm g$};\draw[axis](-.5,2.6)--(-.5,1.4)node[left]{$\bm g$};\node at(1.5,-.5){引力场中的小舱};
\begin{scope}[shift={(6,0)}]
\draw[world](0,0)rectangle(3,3);\fill(1.5,1.5)circle(3pt);\draw[axis](1.5,1.4)--(1.5,.5)node[right]{$-m\bm a$};\draw[axis](3.6,1.4)--(3.6,2.6)node[right]{$\bm a$};\node at(1.5,-.5){远离引力场的加速舱};
\end{scope}
\end{tikzpicture}}
局部条件不可省略。跨越较大区域时，引力方向和大小可能随位置变化，产生潮汐效应；一个整体平动加速的参考系不能消去这些变化。

\subsection{轨道、光线与时钟}

广义相对论用时空几何描述引力，可以计算行星轨道的相对论修正、掠过天体的光线偏折，以及不同高度时钟的速率差。这里先给出弱场、低速情形下的时钟关系。

对静止于弱引力场中的时钟，若 $\Phi$ 为引力势且 $|\Phi|/c^2\ll1$，则
\[
\dd\tau\simeq\left(1+\frac{\Phi}{c^2}\right)\dd t.
\]
比较图\ref{fig:hr-gr2}中的上下两处，$\Phi_2>\Phi_1$，高处时钟相对于低处走得较快。低处发出的光到达高处后，由高处本地时钟测得的频率降低，近似满足
\[
\frac{\nu_2-\nu_1}{\nu_1}\simeq-\frac{\Phi_2-\Phi_1}{c^2}.
\]
\fig{fig:hr-gr2}{静态弱引力场中的时钟比较与向上传播的光}{
\begin{tikzpicture}[scale=.95]
\draw[world](-.5,0)--(6.5,0);\draw[world](-.5,2.5)--(6.5,2.5);\draw[axis](-1.1,2)--(-1.1,.5)node[midway,left]{$\bm g$};
\draw[world](0,0)circle(.3);\node[above right]at(.5,0){$\Phi_1,\ \nu_1$};\draw[world](0,2.5)circle(.3);\node[above right]at(.5,2.5){$\Phi_2>\Phi_1,\ \nu_2<\nu_1$};
\draw[light](4.7,.2)--(4.7,2.3)node[midway,right]{光向上传播};
\end{tikzpicture}}
引力势决定上述时钟比较，不能仅凭某一点的引力场强大小判断时钟快慢。狭义相对论的运动学效应与引力效应在实际比较中还可能同时出现。
例如在地面附近取较小高度差 $h$，近似有 $\Phi_2-\Phi_1\simeq gh$，于是
\[
\frac{\dd\tau_2}{\dd\tau_1}\simeq1+\frac{gh}{c^2},\qquad
\frac{\nu_2-\nu_1}{\nu_1}\simeq-\frac{gh}{c^2}.
\]
频率变化的符号也可从波峰计数理解。静态场中，连续两波峰到达上下两处的坐标时间间隔相同，记为 $\dd t$；本地测得的周期却分别为 $\dd\tau_1,\dd\tau_2$。因而
\[
 \frac{\nu_2}{\nu_1}=\frac{\dd\tau_1}{\dd\tau_2}
 \simeq\frac{1+\Phi_1/c^2}{1+\Phi_2/c^2}
 \simeq1-\frac{\Phi_2-\Phi_1}{c^2}.
\]
频率比较必须使用各处本地钟，而不能把两个位置的固有时间先当成相同。

图\ref{fig:hr-gr2}的上方时钟每一单位坐标时间走得稍快，向上传播的光在上方用本地钟计量时频率稍低；两个符号描述的是同一次比较的不同测量量。这里要求静态弱场、上下两处相对静止，且高度范围内 $g$ 的变化可以忽略。

引力场的扰动也能以波的形式传播。引力波在局部表现为横向两个方向上交替的潮汐伸缩，可通过比较互相垂直的光程变化进行测量。理解这类现象需要把本讲的时钟、光信号和事件观点推广到弯曲时空。

\section{例题与分析}

\begin{question}
\begin{example}[炸弹爆炸事件的坐标]\label{ex:hr1}
在地面系 $S$ 中，炸弹于 $x=1.0\times10^6\,\mathrm m,t=0.020\,\mathrm s$ 爆炸。飞船系 $S'$ 沿 $+x$ 以 $V=0.75c$ 运动，两系原点在 $t=t'=0$ 重合。求爆炸事件在飞船系中的地点和时刻，并与伽利略变换比较，见图\ref{fig:hr1q}。
\end{example}
\fig{fig:hr1q}{地面系中爆炸事件 $A$ 与飞船原点的世界线。两轴均以 $10^6\,\mathrm m$ 为单位，等比例。}{
\begin{tikzpicture}[scale=.65]
\draw[axis](-.4,0)--(5.5,0)node[right]{$x/(10^6\,\mathrm m)$};\draw[axis](0,-.4)--(0,7)node[above]{$ct/(10^6\,\mathrm m)$};
\draw[world](0,0)--(4.95,6.6)node[right]{$x=0.75ct$};\fill(1,6)circle(2pt)node[above right]{$A(1,6)$};\draw[aux](1,0)--(1,6)--(0,6);\node[below left]at(0,0){$O$};
\end{tikzpicture}}
\end{question}

\begin{solution}\textbf{统一时空单位并求洛伦兹因子}\par\nobreak
图\ref{fig:hr1q}中时间坐标已乘以 $c$，所以 $ct_A=6.0\times10^6\,\mathrm m$。先求
\[
\gamma=\frac1{\sqrt{1-0.75^2}}=\frac4{\sqrt7}.
\]


\par\medskip\Needspace{4\baselineskip}\noindent\textbf{变换同一爆炸事件的两个坐标}\par\nobreak
把同一个事件的空间、时间坐标一起变换：
\[
x'_A=\frac4{\sqrt7}(1.0-0.75\times6.0)10^6\,\mathrm m
=\boxed{-5.2915\times10^6\,\mathrm m},
\]
\[
t'_A=\frac4{\sqrt7}\left(0.020-\frac{0.75\times10^6}{3.00\times10^8}\right)\mathrm s
=\boxed{0.026458\,\mathrm s}.
\]
\fig{fig:hr1s}{飞船系中的同一事件。负横坐标表示爆炸发生在飞船原点后方，时间仍为正。}{
\begin{tikzpicture}[scale=.55]
\draw[axis](-6.4,0)--(1,0)node[right]{$x'/(10^6\,\mathrm m)$};\draw[axis](0,-.3)--(0,9)node[above]{$ct'/(10^6\,\mathrm m)$};
\fill(-5.291503,7.937254)circle(2pt)node[above left]{$A$};\draw[aux](-5.291503,0)--(-5.291503,7.937254)--(0,7.937254);
\node[below]at(-5.291503,0){$-5.2915$};\node[right]at(0,7.937254){$7.9373$};
\end{tikzpicture}}



\par\medskip\Needspace{4\baselineskip}\noindent\textbf{用位置趋势和不变间隔复核}\par\nobreak
地面爆炸时，飞船原点已在 $Vt=4.5\times10^6\,\mathrm m$ 处，所以图\ref{fig:hr1s}中事件在后方是合理的。间隔也可核对：
\[
(ct_A)^2-x_A^2=35\times10^{12}\,\mathrm{m^2}
=(ct'_A)^2-x_A'^2.
\]


\par\medskip\Needspace{4\baselineskip}\noindent\textbf{比较经典结果并区别信号到达时间}\par\nobreak
若用伽利略变换，则 $x'_{\rm G}=x-Vt=-3.5\times10^6\,\mathrm m$，$t'_{\rm G}=0.020\,\mathrm s$。在 $0.75c$ 下不能忽略时间变换及洛伦兹因子。
图\ref{fig:hr1s}的 $t'_A$ 由飞船系同步钟网定义，表示爆炸发生的坐标时间；光信号随后到达原点，还需传播一段时间。

又因 $O,A$ 为类时间隔且 $t_A>0$，所有惯性系都把 $A$ 排在 $O$ 之后。$x'_A<0$ 表示爆炸地点在后方，与这个时间先后关系相容。
\end{solution}

\begin{question}
\begin{example}[隧道口的闪电能否击中车尾]\label{ex:hr5}
隧道 $BA$ 与列车 $B'A'$ 的固有长度都为 $L_0$。列车以 $0<V<c$ 沿 $B\to A$ 方向运动。在地面系中，车头 $A'$ 恰到出口 $A$ 的同时，闪电击中入口 $B$。问闪电是否同时击中车尾 $B'$？见图\ref{fig:hr5q}。将闪电击中视为入口处的一次局部事件。
\end{example}
\fig{fig:hr5q}{所画快照属于地面系：隧道长 $L_0$，列车同时测长为 $L_0/\gamma$，车尾已进入入口右侧。}{
\begin{tikzpicture}
\draw[thick](0,1.2)--(5,1.2) (0,0)--(0,1.2) (5,0)--(5,1.2);\draw[very thick](1,0)--(5,0);
\node[below]at(0,0){$B$};\node[below]at(1,0){$B'$};\node[below]at(5,0){$A,A'$};\draw[axis](1.5,.5)--(3.5,.5)node[midway,above]{$V$};\draw[dim](0,1.7)--(5,1.7)node[midway,above]{$L_0$};
\draw[thick](-.2,2.3)--(.15,1.9)--(-.1,1.5)--(0,.9);\node[left]at(-.25,1.8){闪电};
\end{tikzpicture}}
\end{question}

\begin{solution}\textbf{在地面同时面上确定车尾位置}\par\nobreak
取入口 $B$ 为 $x=0$，出口 $A$ 为 $x=L_0$，并把题设地面同时的两事件定为 $t=0$。此时车头在 $x=L_0$，车长却是 $L_0/\gamma$，所以车尾在
\[
x_{B'}(0)=L_0-L_0/\gamma>0.
\]


\par\medskip\Needspace{4\baselineskip}\noindent\textbf{以局部事件是否重合判断击中}\par\nobreak
闪电事件在 $(x,t)=(0,0)$，车尾事件在 $(L_0-L_0/\gamma,0)$，二者不同地。因此 $\boxed{\text{闪电不会击中车尾}}$。
\fig{fig:hr5s}{地面时空图。闪电事件 $E$ 在入口世界线上，车尾世界线在其右方；车头抵出口事件为 $F$。两轴等比例，图示速度取 $0.6c$。}{
\begin{tikzpicture}[scale=.65]
\draw[axis](-.4,0)--(8.4,0)node[right]{$x$};\draw[axis](0,-.4)--(0,5.4)node[above]{$ct$};\draw[world](5,0)--(5,5);
\draw[world](1,0)--(4,5);\draw[world](5,0)--(8,5);\fill(0,0)circle(2pt)node[below left]{$E$};\fill(5,0)circle(2pt)node[below right]{$F$};
\node[left]at(0,4.4){入口};\node[above right]at(5,5){出口};\node[above left]at(4,5){车尾};\node[above]at(8,5){车头};\draw[aux](0,0)--(5,3)node[right]{$t'=0$};
\end{tikzpicture}}



\par\medskip\Needspace{4\baselineskip}\noindent\textbf{将同一对事件变换到车系}\par\nobreak
“击中”要求两个物体在同一事件相遇，这种重合关系在所有参考系中一致。在车系中，$E,F$ 不再同时。以 $E$ 为两系共同原点，时间变换给出
\[
t'_E=0,\qquad t'_F=-\gamma\frac{VL_0}{c^2}<0.
\]
在车系的闪电时刻 $t'=0$，还可以直接列位置。



\par\medskip\Needspace{4\baselineskip}\noindent\textbf{重选车系同时面并画出列车位置}\par\nobreak
入口在 $x'_B=0$，出口在 $x'_A=L_0/\gamma$，而车尾、车头分别固定于
\[
x'_{B'}=(\gamma-1)L_0,\qquad x'_{A'}=\gamma L_0.
\]
因此车尾与入口仍有正的间隔，结论与地面系一致。图\ref{fig:hr5ship}采用与前图相同的 $V=0.6c$，把这组同时位置画出；此时车头早已超过出口。
\fig{fig:hr5ship}{车系 $t'=0$ 的同时快照。隧道以 $-V$ 运动，列车静止；本图的出口位置与地面同时快照不同。}{
\begin{tikzpicture}[scale=.95]
\draw[world](0,0)--(0,1.2)--(4,1.2)--(4,0);\draw[line width=1.3pt](1.25,0)--(6.25,0);
\node[below]at(0,0){$B$};\node[below]at(1.25,0){$B'$};\node[above right]at(4,.05){$A$};\node[below]at(6.25,0){$A'$};
\draw[dim](0,1.7)--(4,1.7)node[midway,above]{$L_0/\gamma$};\draw[dim](1.25,-.7)--(6.25,-.7)node[midway,below]{$L_0$};
\draw[axis](3.3,.65)--(1.8,.65)node[midway,above]{$-V$};\draw[world](-.25,1.6)--(.1,1.15)--(-.15,.7)--(0,.15);\node[left]at(-.25,1.6){闪电};
\end{tikzpicture}}
图\ref{fig:hr5s}的斜同时线说明了两组快照为何不同：车系中，车头抵出口发生在闪电之前。



\par\medskip\Needspace{4\baselineskip}\noindent\textbf{区别局部击中和随后光信号传播}\par\nobreak
若另问闪电发出的光后来能否追上车尾，则需画出光信号世界线，寻找它与车尾世界线的交点。
\end{solution}

\begin{question}
\begin{example}[高空产生的介子能否到达地面]\label{ex:hr7}
一个介子在离地面 $D=6000\,\mathrm m$ 的高空产生，以 $v=0.998c$ 飞向地面。假定平均固有寿命为 $\tau_0=2.0\times10^{-6}\,\mathrm s$，按指数规律衰变，忽略运动中的能量损失。（1）地面能否观测到到达的介子？（2）与介子共动的观察者如何判断？见图\ref{fig:hr7q}。
\end{example}
\fig{fig:hr7q}{产生事件为 $A$，到达地面为 $B$。平均寿命是统计参数，并非每个介子都在该时刻消失。}{
\begin{tikzpicture}
\draw[thick](-2,0)--(2,0);\fill(0,3.5)circle(2pt)node[above]{$A$};\draw[axis](0,3.2)--(0,1.6)node[midway,right]{$0.998c$};\node[below]at(0,0){$B$，地面};\draw[dim](-1,0)--(-1,3.5)node[midway,left]{$6000\,\mathrm m$};
\end{tikzpicture}}
\end{question}

\begin{solution}\textbf{选取粒子产生和到达的同一对事件}\par\nobreak
图\ref{fig:hr7s}中，$A,B$ 位于同一粒子的世界线上，实验室飞行时间和粒子固有时间分别为 $D/v$、$D/(\gamma v)$。先求
\[
\gamma=\frac1{\sqrt{1-0.998^2}}\simeq15.8193,
\quad \tau_{\text{地}}=\gamma\tau_0\simeq31.639\,\mu\mathrm s.
\]


\par\medskip\Needspace{4\baselineskip}\noindent\textbf{在地面系计算平均距离与存活比例}\par\nobreak
地面系平均传播距离为 $v\gamma\tau_0\simeq9.47\,\mathrm{km}$，大于 $6.00\,\mathrm{km}$。更确切地，到达地面的存活比例为
\[
\boxed{\frac N{N_0}=\exp\left[-\frac D{\gamma v\tau_0}\right]\simeq0.531.}
\]
即约有 $53.1\%$ 的介子能够到达地面；个别粒子仍可能在途中衰变。
\fig{fig:hr7s}{把沿飞行方向的坐标改为 $x$。探测面静止，$AB$ 为粒子世界线；两轴以相同长度单位绘制。}{
\begin{tikzpicture}[scale=.65]
\draw[axis](-.3,0)--(7.2,0)node[right]{$x/\mathrm{km}$};\draw[axis](0,-.3)--(0,7)node[above]{$ct/\mathrm{km}$};\draw[world](6,0)--(6,6.7);\draw[world](0,0)--(6,6.012024);\fill(6,6.012024)circle(2pt)node[right]{$B$};\node[below left]at(0,0){$A$};\draw[aux](0,6.012024)--(6,6.012024);\node[above]at(6,6.7){地面};\node[rotate=45,above]at(2.8,2.8){$\tau_{AB}=D/(\gamma v)$};
\end{tikzpicture}}



\par\medskip\Needspace{4\baselineskip}\noindent\textbf{在粒子系重选地面同时位置}\par\nobreak
介子系中，地面与产生位置在地面系的静止间距收缩为 $D'=D/\gamma\simeq379.3\,\mathrm m$。



\par\medskip\Needspace{4\baselineskip}\noindent\textbf{用粒子固有时间复核衰变统计}\par\nobreak
地面以 $0.998c$ 接近，所需时间
\[
t'_{AB}=\frac{D'}v=\frac D{\gamma v}\simeq1.267\,\mu\mathrm s.
\]
将它代入同一指数衰变式，仍得约 $0.531$。两系选取的是相同的产生、到达事件，因而关于能否到达的结论一致。
\fig{fig:hr7rest}{介子系中的两幅同时快照。介子始终在 $x'=0$，地面从 $D/\gamma$ 处以 $-v$ 接近。}{
\begin{tikzpicture}[scale=.9]
\draw[axis](-.4,1.8)--(5.2,1.8)node[right]{$x'$};\fill(0,1.8)circle(2pt)node[below]{$A$};\draw[world](4,1.4)--(4,2.2)node[right]{地面};\draw[dim](0,2.6)--(4,2.6)node[midway,above]{$D'=D/\gamma$};\draw[axis](3.6,1.05)--(2,1.05)node[midway,below]{$-v$};\node[left]at(-.6,1.8){$t'=0$};
\draw[axis](-.4,-.5)--(5.2,-.5)node[right]{$x'$};\draw[world](0,-.9)--(0,-.1);\fill(0,-.5)circle(2pt)node[below right]{$B$};\node[left]at(-.6,-.5){$t'=D' /v$};\node[above right]at(1,-.5){地面到达介子处};
\end{tikzpicture}}
图\ref{fig:hr7rest}中，地面到达所需的时间为 $1.267\,\mu\mathrm s$，介子的平均固有寿命为 $2.0\,\mu\mathrm s$。平均寿命描述一群粒子的衰变统计，单个粒子没有统一的消失时刻。



\par\medskip\Needspace{4\baselineskip}\noindent\textbf{检查飞行时间和平均寿命所属的系}\par\nobreak
计算存活比例时，飞行时间与寿命应取自同一参考系；介子系用收缩后的距离 $D/\gamma$，地面系用膨胀后的平均寿命 $\gamma\tau_0$。
\end{solution}

\begin{question}
\begin{example}[反向飞行的两艘飞船]\label{ex:hr2}
地面测得飞船 $a,b$ 沿同一直线分别以 $+0.9c,-0.9c$ 飞行。求 $a$ 相对 $b$ 的速度大小，见图\ref{fig:hr2q}。
\end{example}
\fig{fig:hr2q}{两速度都给在地面系中，正方向向右。}{
\begin{tikzpicture}
\draw[axis](-4,0)--(4,0)node[right]{$x$};\fill(-2.8,0)circle(2pt)node[below]{$a$};\fill(2.8,0)circle(2pt)node[below]{$b$};
\draw[axis](-3,.7)--(-.7,.7)node[midway,above]{$+0.9c$};\draw[axis](3,.7)--(.7,.7)node[midway,above]{$-0.9c$};
\end{tikzpicture}}
\end{question}

\begin{solution}\textbf{取另一飞船静止系并确定变换符号}\par\nobreak
以 $b$ 的静止系为 $S'$，该系相对地面速度 $V=-0.9c$。把 $a$ 的地面速度 $u_a=+0.9c$ 代入从地面到 $S'$ 的速度变换，
\[
u'_a=\frac{u_a-V}{1-u_aV/c^2}
=\frac{0.9c+0.9c}{1+0.81}
=\boxed{\frac{180}{181}c\simeq0.994475c}.
\]
\fig{fig:hr2s}{在 $b$ 系中，$b$ 静止，$a$ 沿正方向接近。分母中的时间变换使相对速度保持小于 $c$。}{
\begin{tikzpicture}
\draw[axis](-4,0)--(2,0)node[right]{$x'$};\fill(0,0)circle(2pt)node[below]{$b$，静止};\fill(-3.4,0)circle(2pt)node[below]{$a$};
\draw[axis](-3.4,.7)--(-.7,.7)node[midway,above]{$u'_a=180c/181$};
\end{tikzpicture}}



\par\medskip\Needspace{4\baselineskip}\noindent\textbf{区分地面间距减少率与相对速度}\par\nobreak
图\ref{fig:hr2q}中的地面间距以 $1.8c$ 的速率减少，这是按地面时间计算的间距变化率。图\ref{fig:hr2s}所求速度使用 $b$ 系的尺和同步钟；变换相邻事件时，空间与时间间隔都会改变。由它们的比值得到 $u'_a=180c/181<c$。


\par\medskip\Needspace{4\baselineskip}\noindent\textbf{逆变换复核参考系的速度符号}\par\nobreak
方向也可用逆变换检查。把 $u'_a=180c/181$ 与 $V=-0.9c$ 代入
\[
u_a=\frac{u'_a+V}{1+Vu'_a/c^2},
\]
恰好还原为 $+0.9c$。若把参考系速度误取为 $+0.9c$，便会将“变到 $b$ 系”错当成“变到 $a$ 系”，这是本题符号选择的主要区别。



\par\medskip\Needspace{4\baselineskip}\noindent\textbf{从同一飞船的相邻事件解释时间分母}\par\nobreak
也可把图\ref{fig:hr2s}所求速度直接写成事件间隔之比。取飞船 \(a\) 上相邻两事件，在地面系有 \(\Delta x=0.9c\Delta t\)。变到 \(b\) 系时，
\[
 \Delta x'=\gamma_b(1.8c\Delta t),\qquad
 \Delta t'=\gamma_b(1.81\Delta t).
\]
二者相除，公共因子 \(\gamma_b\) 消去，得到 \(u'_a=1.8c/1.81\)。时间分母的变化就是不能直接把两个 \(0.9c\) 相加的原因。
\end{solution}

\begin{question}
\begin{example}[飞船向前发射导弹]\label{ex:hr4}
飞船以 $0.9c$ 远离地球，又以相对飞船 $0.9c$ 的速率向前发射导弹。求导弹相对地球的速率，见图\ref{fig:hr4q}。
\end{example}
\fig{fig:hr4q}{船系中的发射速度为 $u'=0.9c$，飞船相对地球速度为 $V=0.9c$。}{
\begin{tikzpicture}
\draw(0,-.5)rectangle(3.5,.5);\node at(1.75,0){飞船 $S'$};\draw[axis](3.5,0)--(5.8,0)node[right]{$u'=0.9c$};\draw[axis](.7,1.2)--(2.8,1.2)node[right]{$V=0.9c$};
\end{tikzpicture}}
\end{question}

\begin{solution}\textbf{由船系发射速度合成为地球系速度}\par\nobreak
图\ref{fig:hr4q}的两个速度属于不同的参考系，应由逆速度变换合成：
\[
u=\frac{u'+V}{1+Vu'/c^2}
=\frac{1.8}{1.81}c
=\boxed{\frac{180}{181}c\simeq0.994475c}.
\]
\fig{fig:hr4s}{地球系中，导弹比飞船快，但仍比真空光慢。三条箭头按速度比例绘制。}{
\begin{tikzpicture}
\draw[axis](0,1)--(5.96685,1)node[right]{$u=180c/181$};\draw[axis](0,0)--(5.4,0)node[right]{$V=0.9c$};\draw[axis](0,-1)--(6,-1)node[right]{$c$};\node[left]at(0,1){导弹};\node[left]at(0,0){飞船};\node[left]at(0,-1){光};
\end{tikzpicture}}



\par\medskip\Needspace{4\baselineskip}\noindent\textbf{检查速度差与光速上限}\par\nobreak
图\ref{fig:hr4s}中，地球系速度差 $u-V\simeq0.094475c$ 不是导弹的船系速度。两系赋予相邻导弹事件的时间间隔不同，须把空间和时间一起变换。
还可把合成速度与光速的差写成
\[
c-u=\frac{(c-u')(c-V)}{c(1+Vu'/c^2)}.
\]
本题右端为 $0.01c/1.81=c/181>0$，与图\ref{fig:hr4s}导弹箭头略短于光箭头一致。



\par\medskip\Needspace{4\baselineskip}\noindent\textbf{以导弹相邻事件推导所用比值}\par\nobreak
为说明这里使用的是逆变换，取导弹在船系的两相邻事件，满足 \(\Delta x'=u'\Delta t'\)。地球系测得
\[
 \Delta x=\gamma_V(u'+V)\Delta t',\qquad
 \Delta t=\gamma_V(1+Vu'/c^2)\Delta t'.
\]
与图\ref{fig:hr4s}向右的传播方向对应，两式取比正好得到前述合成式。反向计算 $(u-V)/(1-uV/c^2)$ 恰好得到船系的 $0.9c$。低速极限下，当 \(Vu'/c^2\ll1\) 时才近似回到 \(u\simeq u'+V\)；本题该无量纲量为 \(0.81\)，显然不能忽略。
\end{solution}

\begin{question}
\begin{example}[飞船与彗星的相撞时间]\label{ex:hr6}
飞船和彗星相对地面分别以 $+0.6c,-0.8c$ 相向运动。地面在当前时刻判断，再过 $5\,\mathrm s$ 两者相撞。求：（1）飞船系中彗星的速度；（2）从飞船当前事件到相撞，船钟还经过多少时间？见图\ref{fig:hr6q}。
\end{example}
\fig{fig:hr6q}{地面当前快照。两者间距由地面接近率与剩余时间确定。}{
\begin{tikzpicture}
\draw[axis](-.5,0)--(7.5,0)node[right]{$x$};\fill(0,0)circle(2pt)node[below]{飞船};\fill(7,0)circle(2pt)node[below]{彗星};\draw[axis](.3,.7)--(2.3,.7)node[midway,above]{$0.6c$};\draw[axis](6.7,.7)--(4.7,.7)node[midway,above]{$0.8c$};\draw[dim](0,-.7)--(7,-.7)node[midway,below]{$D=(0.6+0.8)c\times5\,\mathrm s$};
\end{tikzpicture}}
\end{question}

\begin{solution}\textbf{在地面系写出两条运动世界线}\par\nobreak
取飞船当前事件为 $O$，令 $x_O=t_O=0$。图\ref{fig:hr6s}中两条世界线为
\[
x_{\text{船}}=0.6ct,\qquad x_{\text{彗}}=7c\,\mathrm s-0.8ct.
\]


\par\medskip\Needspace{4\baselineskip}\noindent\textbf{把彗星速度变到飞船系}\par\nobreak
彗星速度变到飞船系，
\[
\boxed{u'_{\text{彗}}=\frac{-0.8c-0.6c}{1+0.48}
=-\frac{35}{37}c\simeq-0.945946c.}
\]


\par\medskip\Needspace{4\baselineskip}\noindent\textbf{在飞船世界线上计算船钟走时}\par\nobreak
相撞事件 $C$ 在地面系有 $t_C=5\,\mathrm s,x_C=3c\,\mathrm s$。由于 $O,C$ 都在飞船世界线上，船钟经历
\[
\boxed{\Delta\tau_{OC}=\frac5{1.25}\,\mathrm s=4.00\,\mathrm s.}
\]
\fig{fig:hr6s}{横轴 $x/c$、纵轴 $t$ 均以秒计且等比例。虚线 $OE_*$ 在飞船系同时，$E_*$ 不是地面 $t=0$ 时的彗星事件。}{
\begin{tikzpicture}[scale=.65]
\draw[axis](-.4,0)--(8,0)node[right]{$x/c\ (\mathrm s)$};\draw[axis](0,-.3)--(0,6.5)node[above]{$t\ (\mathrm s)$};
\draw[world](0,0)--(3.6,6);\draw[world](7,0)--(2.2,6);\fill(3,5)circle(2pt)node[above right]{$C(3,5)$};\draw[aux](0,0)--(4.72973,2.83784);\fill(4.72973,2.83784)circle(2pt)node[right]{$E_*$};\node[below left]at(0,0){$O$};\node[below]at(7,0){彗星初始事件};\node[above,rotate=30.964]at(2,1.2){$t'=0$};
\end{tikzpicture}}



\par\medskip\Needspace{4\baselineskip}\noindent\textbf{重新求船系同时的彗星位置并复核时间}\par\nobreak
用时间变换也得 $\Delta t'=1.25(5-0.6\times3)\,\mathrm s=4\,\mathrm s$。若改用“船系初始距离除以相对速度”，须先按图\ref{fig:hr6s}的船系同时线取彗星位置：$t=0.6x/c$ 与其世界线相交，得
\[
\frac{x}{c}=7\,\mathrm s-0.8t
=7\,\mathrm s-0.48\frac{x}{c},\qquad
1.48\frac{x}{c}=7\,\mathrm s.
\]
所以
\[
\frac{x_{E_*}}c=\frac{175}{37}\,\mathrm s,\quad
t_{E_*}=\frac{105}{37}\,\mathrm s,\quad
D'=x'_{E_*}=\frac{140}{37}c\,\mathrm s.
\]
于是 $D'/|u'_{\text{彗}}|=4\,\mathrm s$。不能把地面同时距离 $D$ 直接除以船系相对速度。
\end{solution}

\begin{question}
\begin{example}[运动圆孔与固定挡板的遮光时刻]\label{ex:hr9}
车厢以速度 $V$ 沿地面 $+x$ 运动。车内点光源 $S$ 一直发光，正前方固有距离 $l$ 处有半径 $r$ 的圆孔，圆心为 $O_1$。地面 $x=x_A$ 处固定半径 $R>r$ 的不透明圆挡板，圆心为 $O_2$。孔面、挡板均垂直于轴线，$S,O_1,O_2$ 等高，见图\ref{fig:hr9q}。$S$ 经过地面原点时取 $t=t'=0$，此时尚有光从挡板周围通过。求挡板处刚完全截住光束这一事件的车系时刻与地面时刻。忽略衍射；不把挡板到远处屏幕的传播延迟计入。
\end{example}
\fig{fig:hr9q}{在车系中作几何示意：光源和圆孔静止，挡板向左运动，边界光线形成稳定的圆锥。}{
\begin{tikzpicture}[scale=.9]
\draw[axis](-.5,0)--(8.2,0)node[right]{$x'$};\draw[very thick](2,.6)--(2,2.3) (2,-.6)--(2,-2.3);\draw[very thick](6,-1.2)--(6,1.2);
\draw[light](0,0)--(2,.6);\draw[light](2,.6)--(8,2.4);\draw[light](0,0)--(2,-.6);\draw[light](2,-.6)--(8,-2.4);
\fill(0,0)circle(1.8pt)node[left]{$S$};\node[below left]at(2,0){$O_1$};\node[below right]at(6,0){$O_2$};\draw[dim](2.4,0)--(2.4,.6)node[midway,right]{$r$};\draw[dim](6.4,0)--(6.4,1.2)node[midway,right]{$R$};\draw[dim](0,-1.1)--(2,-1.1)node[midway,below]{$l$};\draw[axis](6.4,2)--(4.7,2)node[midway,above]{$-V$};
\end{tikzpicture}}
\end{question}

\begin{solution}\textbf{在车系由相似三角形求遮光位置}\par\nobreak
车系中，圆孔和光源均静止，边界光线由 $S$ 连到孔边缘。因此图\ref{fig:hr9s}中距光源 $X'$ 处的光束半径为 $\rho=rX'/l$。刚完全遮光时 $\rho=R$，故挡板的车系坐标为
\[
\boxed{X'_*=\frac{lR}{r}}.
\]
由于 $R>r$，该位置仍在孔面前方。
\fig{fig:hr9s}{遮光开始时，挡板边缘恰好接触光锥边界。两直角三角形给出 $R/X'_*=r/l$。}{
\begin{tikzpicture}
\draw[axis](-.5,0)--(5.2,0)node[right]{$x'$};\draw[very thick](2,.6)--(2,1.8) (2,-.6)--(2,-1.8);\draw[very thick](4,-1.2)--(4,1.2);
\draw[light](0,0)--(2,.6);\draw[light](2,.6)--(4,1.2);\draw[light](0,0)--(2,-.6);\draw[light](2,-.6)--(4,-1.2);
\node[left]at(0,0){$S$};\node[above]at(4,1.2){$E$，边缘接触};\draw[dim](4.4,0)--(4.4,1.2)node[midway,right]{$R$};\draw[dim](2.3,0)--(2.3,.6)node[midway,right]{$r$};\draw[dim](0,-2.2)--(4,-2.2)node[midway,below]{$X'_*=lR/r$};
\end{tikzpicture}}



\par\medskip\Needspace{4\baselineskip}\noindent\textbf{变换挡板世界线并求车系时刻}\par\nobreak
地面固定挡板的世界线为 $x=x_A$。由逆变换 $x=\gamma(x'+Vt')$，在车系有
\[
x'_A(t')=\frac{x_A}{\gamma}-Vt'.
\]
令它到达 $X'_*$，得到车系坐标时间
\[
\boxed{t'_E=\frac1V\left(\frac{x_A}{\gamma}-\frac{lR}{r}\right).}
\]
这里的 $t'_E$ 是车系同步钟网赋予挡板处遮光事件的坐标时间。

\fig{fig:hr9time}{地面时空图。$x'=X'_*$ 对应一条与车厢平行的直线，与固定挡板世界线相交于遮光事件 $E$。两轴等比例，图示取 $V=0.6c$。}{
\begin{tikzpicture}[scale=.5]
\draw[axis](-.5,0)--(10.2,0)node[right]{$x$};\draw[axis](0,-.4)--(0,10)node[above]{$ct$};\draw[world](0,0)--(5.7,9.5)node[above]{$S$};\draw[world](8,0)--(8,9.5)node[above]{挡板};\draw[aux](3.2,0)--(8.9,9.5);\fill(8,8)circle(2pt)node[right]{$E$};\node[above,rotate=59.036]at(6.5,5.5){$x'=X'_*$};\node[below]at(8,0){$x_A$};\node[below left]at(0,0){$O$};
\end{tikzpicture}}


\par\medskip\Needspace{4\baselineskip}\noindent\textbf{把同一遮光事件变回地面时刻}\par\nobreak
图\ref{fig:hr9time}中遮光事件 $E$是固定挡板世界线与 $x'=X'_* $轨迹的交点，先确定这个共同事件，再变回地面时间；等价地，将 $x'=\gamma(x_A-Vt)$ 代入阈值，得
\[
\boxed{t_E=\frac1V\left(x_A-\frac{lR}{\gamma r}\right).}
\]


\par\medskip\Needspace{4\baselineskip}\noindent\textbf{核对初始条件和两种时间零点}\par\nobreak
题给地面零时刻仍有光从边缘通过，等价于 $\gamma x_A>lR/r$，保证 $t_E>0$。$t'_E$ 不一定为正，因为地面零时刻与车系零时刻的远处同时面不同。若另求屏幕真正变暗的接收时刻，还须给出屏幕位置并加上最后一束光从挡板到屏幕的传播时间。



\par\medskip\Needspace{4\baselineskip}\noindent\textbf{检查固有距离和横向尺寸的使用}\par\nobreak
在几何步骤中，$l$ 是车厢固有距离，图\ref{fig:hr9s}又正是在车系中作图，所以直接用 $l$，不再附加一次长度收缩。挡板的横向半径 $R$ 和圆孔半径 $r$ 也不随纵向变换改变。遮光发生在一圈边缘的同一轴向位置、同一时刻，时空图中的 $E$ 代表其中任一点的轴向坐标。



\par\medskip\Needspace{4\baselineskip}\noindent\textbf{用挡板半径变化检查结论}\par\nobreak
最后检查参数趋势：固定 $V,x_A,l,r$ 时，
\[
\frac{\partial t_E}{\partial R}=-\frac{l}{\gamma rV}<0.
\]
挡板越大，越早完全遮光，符合光锥与挡板相交的图形。反之，若圆孔更大，光束更宽，挡板需更接近光源才会截住全部光。
\end{solution}

\begin{question}
\begin{example}[串列静电加速器]\label{ex:hr8}
串列加速器中，离子源 $S$ 与终端靶 $T$ 接地，中央氮气管 $N$ 的电势为 $U=5.00\times10^6\,\mathrm V$，见图\ref{fig:hr8q}。氢负离子从 $S$ 由静止进入，在 $N$ 中被剥去电子而成为质子。忽略电子质量及能量损失，并假定剥离不改变粒子速度。考虑相对论效应，求质子到达 $T$ 的速率。取 $e=1.60\times10^{-19}\,\mathrm C,m_0=1.673\times10^{-27}\,\mathrm{kg}$。
\end{example}
\fig{fig:hr8q}{粒子先以负电荷从零电势到正高压区，再以正电荷回到零电势；两段均加速。}{
\begin{tikzpicture}
\draw[thick](0,-.6)rectangle(6,.6);\fill[black!12](2.6,-.5)rectangle(3.4,.5);\draw(2.6,-.5)rectangle(3.4,.5);
\draw[axis](5.7,0)--(3.5,0)node[midway,above]{$\mathrm H^-$};\draw[axis](2.5,0)--(.3,0)node[midway,above]{$\mathrm H^+$};
\node[above]at(6,.7){$S$};\node[above]at(0,.7){$T$};\node[below]at(3,-.6){$N$，$+U$};\node[below]at(0,-.6){$0$};\node[below]at(6,-.6){$0$};
\draw(3,.5)--(3,1.5);\draw(2.5,1.5)rectangle(3.5,2.1);\node at(3,1.8){$H$};
\end{tikzpicture}}
\end{question}

\begin{solution}\textbf{按电荷和电势差分别计算两段做功}\par\nobreak
电场做功为 $\Delta K=-q\Delta\varphi$。图\ref{fig:hr8q}的第一段 $S\to N$ 中，$q=-e,\Delta\varphi=U$，所以 $\Delta K_1=eU$。第二段 $N\to T$ 中电荷变成 $+e$，电势降为 $\Delta\varphi=-U$，故仍有 $\Delta K_2=eU$。
\fig{fig:hr8s}{沿粒子行程画电势和动能：电势先升后降，电荷在 $N$ 变号，使动能两段都增加。横轴表示行程次序。}{
\begin{tikzpicture}[scale=.9]
\draw[axis](0,0)--(6.5,0)node[right]{行程};\draw[axis](0,0)--(0,2.5)node[above]{$\varphi$};\draw[thick](0,0)--(3,2)--(6,0);\node[left]at(0,2){$U$};\node[above]at(3,2){$N$};\node[below]at(0,0){$S$};\node[below]at(6,0){$T$};
\begin{scope}[shift={(0,-3.6)}]\draw[axis](0,0)--(6.5,0)node[right]{行程};\draw[axis](0,0)--(0,2.6)node[above]{$K$};\draw[thick](0,0)--(3,1.1)--(6,2.2);\draw[aux](0,1.1)--(3,1.1);\draw[aux](0,2.2)--(6,2.2);\node[left]at(0,1.1){$eU$};\node[left]at(0,2.2){$2eU$};\node[below]at(3,0){剥离};\end{scope}
\end{tikzpicture}}



\par\medskip\Needspace{4\baselineskip}\noindent\textbf{将两段增益相加求终端动能}\par\nobreak
按题设的质量近似，剥离前后的动能可连续处理，终端动能为
\[
K_T=2eU=1.60\times10^{-12}\,\mathrm J=10.0\,\mathrm{MeV}.
\]


\par\medskip\Needspace{4\baselineskip}\noindent\textbf{用相对论动能反求速度}\par\nobreak
由 $K_T=(\gamma-1)m_0c^2$，
\[
\gamma=1+\frac{2eU}{m_0c^2},\qquad
\boxed{v=c\sqrt{1-\left(1+\frac{2eU}{m_0c^2}\right)^{-2}}}
\simeq4.34\times10^7\,\mathrm{m/s}.
\]
图\ref{fig:hr8s}中的两次能量增益都须计入，不能只取一次 $eU$。若要保留电子质量，剥离时速度不变并不等于质子动能不变，须分别使用氢负离子质量和质子质量；本题已明确忽略这一小修正。



\par\medskip\Needspace{4\baselineskip}\noindent\textbf{沿电势图核对两段受力方向}\par\nobreak
两段加速也可从受力方向检查：$S\to N$ 时电势沿运动方向升高，电场指向反方向，负离子受力沿运动方向；$N\to T$ 时电势降低，电场和正离子受力都沿运动方向。因此图\ref{fig:hr8s}的动能在两段均上升。电场做功仍由电荷与电势差决定，相对论修正进入的是最后由 $K_T$ 求 $v$ 的关系。

\end{solution}

\begin{question}
\begin{example}[迎头粘合后的质量]\label{ex:hr3}
两个粒子的不变质量均为 $m$，以大小相同、方向相反的速度 $0.6c$ 迎头碰撞，随后粘成一个复合粒子。若没有能量向外散失，求复合粒子的不变质量 $M$，见图\ref{fig:hr3q}。
\end{example}
\fig{fig:hr3q}{对称的碰前状态。两粒子动量等大反向，总动量为零。}{
\begin{tikzpicture}
\fill(-2,0)circle(3pt)node[below]{$m$};\fill(2,0)circle(3pt)node[below]{$m$};\draw[axis](-2,.6)--(-.5,.6)node[midway,above]{$0.6c$};\draw[axis](2,.6)--(.5,.6)node[midway,above]{$0.6c$};\draw[axisline](-3,0)--(3,0);
\end{tikzpicture}}
\end{question}

\begin{solution}\textbf{先由动量守恒判断复合粒子静止}\par\nobreak
图\ref{fig:hr3s}中两条动量箭头相互抵消，碰前总动量为零。碰后只有一个复合粒子，故其动量也为零，即它在此参考系静止。



\par\medskip\Needspace{4\baselineskip}\noindent\textbf{用总能量守恒求不变质量}\par\nobreak
每个入射粒子的洛伦兹因子为 $\gamma=1/\sqrt{1-0.6^2}=1.25$，总能量为 $2\gamma mc^2$。能量不向外流失，所以
\[
Mc^2=2\gamma mc^2,\qquad \boxed{M=2.50m}.
\]
\fig{fig:hr3s}{上方为碰前的两粒子，下面为静止复合粒子。总动量保持为零，碰前动能转化为复合粒子的内能。}{
\begin{tikzpicture}
\draw[axis](-3,1)--(-1,1)node[midway,above]{$\bm p$};\draw[axis](3,1)--(1,1)node[midway,above]{$-\bm p$};\node[below]at(-2,1){$E=1.25mc^2$};\node[below]at(2,1){$E=1.25mc^2$};
\draw[axis](0,.6)--(0,-.6);\fill(0,-1.1)circle(4pt)node[right]{$M=2.5m,\quad u=0$};\node[below]at(0,-1.6){$E_{\text{后}}=Mc^2=2.5mc^2$};
\end{tikzpicture}}



\par\medskip\Needspace{4\baselineskip}\noindent\textbf{解释增加的静能与内部能量}\par\nobreak
两粒子碰前的静能和为 $2mc^2$，总动能为 $0.5mc^2$。这些相对运动能在粘合后留作内能，使复合粒子的不变质量增加 $0.5m$。若有辐射带走能量和动量，就必须把辐射也列入守恒式，不能沿用本题的 $M$。



\par\medskip\Needspace{4\baselineskip}\noindent\textbf{以系统不变量和一般速率复核}\par\nobreak
也可直接对图\ref{fig:hr3s}的整个系统使用不变量：
\[
M^2c^4=(2\gamma mc^2)^2-c^2|\bm p+(-\bm p)|^2
=4\gamma^2m^2c^4.
\]
取正根得到 $M=2\gamma m$。这样先消去总动量，再求静质量，便不必为碰后的复合粒子另猜一个速度。将 $0.6c$ 换成任意等值反向速率 $v$，仍有 $M=2m/\sqrt{1-v^2/c^2}$。
\end{solution}

\begin{question}
\begin{example}[恒力作用下的高速粒子]\label{ex:hr11}
一粒子不变质量为 $m_0$，起初静止，受地面系中恒力 $F$ 沿 $+x$ 加速，$F/m_0=0.5c/\mathrm s$。求 $t=0,0.1,0.6,1.0\,\mathrm s$ 时的速度和动能，见图\ref{fig:hr11q}。
\end{example}
\fig{fig:hr11q}{恒定的是实验室力 $F=\dd p/\dd t$，不是粒子的加速度。}{
\begin{tikzpicture}
\draw[axis](-.5,0)--(6,0)node[right]{$x$};\fill(.5,0)circle(3pt)node[below]{$m_0$};\draw[axis](.5,.7)--(3,.7)node[midway,above]{$F$，恒定};\node[above]at(.5,1.4){$p(0)=0$};
\end{tikzpicture}}
\end{question}

\begin{solution}\textbf{对恒力积分得到随时间增长的动量}\par\nobreak
图\ref{fig:hr11q}中的 $F$ 是实验室系恒力。把 $\dd p/\dd t=F$ 写为
$\dd p=F\,\dd t$，并使用初值 $p(0)=0$：
\[
 \int_0^{p(t)}\dd p'=\int_0^tF\,\dd t',
 \qquad [p']_0^{p(t)}=F[t']_0^t,\qquad p(t)=Ft.
\]

\par\medskip\Needspace{4\baselineskip}\noindent\textbf{以无量纲动量表达速度和动能}\par\nobreak
令 $q=p/(m_0c)=Ft/(m_0c)$，本题中 $q=0.5t/\mathrm s$。由式\eqref{eq:hrem}，
\[
E=m_0c^2\sqrt{1+q^2},\qquad
\boxed{\frac uc=\frac q{\sqrt{1+q^2}},\quad
\frac K{m_0c^2}=\sqrt{1+q^2}-1.}
\]
\fig{fig:hr11s}{实线为相对论结果，灰虚线为经典近似；$A,B,C,D$ 对应题设四个时刻。}{
\begin{tikzpicture}
\begin{scope}[x=3.6cm,y=4cm]
\draw[axis](0,0)--(1.15,0)node[right]{$t/\mathrm s$};\draw[axis](0,0)--(0,.62)node[above]{$u/c$};\draw[world,domain=0:1.05,samples=60]plot(\x,{.5*\x/sqrt(1+.25*\x*\x)});\draw[gray,dashed](0,0)--(1.05,.525);
\foreach \t/\name in {.6/C,1/D}{\pgfmathsetmacro{\vv}{.5*\t/sqrt(1+.25*\t*\t)}\fill(\t,\vv)circle(1.4pt);\node[above left]at(\t,\vv){$\name$};}
\fill(0,0)circle(1.4pt)node[below left]{$A$};\fill(.1,.0499376)circle(1.4pt);\draw[aux](.1,.0499376)--(.2,.15)node[above]{$B$};
\end{scope}
\begin{scope}[shift={(6.2,0)},x=3.6cm,y=15cm]
\draw[axis](0,0)--(1.15,0)node[right]{$t/\mathrm s$};\draw[axis](0,0)--(0,.17)node[above]{$K/(m_0c^2)$};\draw[world,domain=0:1.05,samples=60]plot(\x,{sqrt(1+.25*\x*\x)-1});\draw[gray,dashed,domain=0:1.05,samples=60]plot(\x,{.125*(\x)^2});
\foreach \t/\name in {.6/C,1/D}{\pgfmathsetmacro{\kk}{sqrt(1+.25*\t*\t)-1}\fill(\t,\kk)circle(1.4pt);\node[above left]at(\t,\kk){$\name$};}
\fill(0,0)circle(1.4pt)node[below left]{$A$};\fill(.1,.00124922)circle(1.4pt);\draw[aux](.1,.00124922)--(.22,.045)node[above]{$B$};
\end{scope}
\end{tikzpicture}}



\par\medskip\Needspace{4\baselineskip}\noindent\textbf{把四个时刻对应到曲线和数值表}\par\nobreak
逐项代入，得
\[
\begin{array}{c|c|c|c}
\text{图中点}&t/\mathrm s&u/c&K/(m_0c^2)\\\hline
A&0&0&0\\
B&0.1&0.049938&0.0012492\\
C&0.6&0.287348&0.0440307\\
D&1.0&0.447214&0.118034
\end{array}
\]


\par\medskip\Needspace{4\baselineskip}\noindent\textbf{求加速度并检查低速极限}\par\nobreak
图\ref{fig:hr11s}的每一点都由同一个动量 $Ft$ 决定。由 $u=cq(1+q^2)^{-1/2}$、$\dd q/\dd t=F/(m_0c)$，
先对 $q$ 使用乘积法则，再作时间的链式求导：
\begin{align*}
 \frac{\dd u}{\dd q}
 &=c\left[(1+q^2)^{-1/2}-q^2(1+q^2)^{-3/2}\right]
 =\frac{c}{(1+q^2)^{3/2}},\\
 \frac{\dd u}{\dd t}
 &=\frac{\dd u}{\dd q}\frac{\dd q}{\dd t}
 =\frac{F/m_0}{(1+q^2)^{3/2}}.
\end{align*}
可见加速度随时间减小，而 $u$ 始终小于 $c$。只有 $q\ll1$ 时才有 $u\simeq Ft/m_0$ 和 $K\simeq F^2t^2/(2m_0)$。


\par\medskip\Needspace{4\baselineskip}\noindent\textbf{由位移积分验证功与动能一致}\par\nobreak
还可用功检查表中的动能。取 $x(0)=0$，将速度对时间积分，并以
$q'=Ft'/(m_0c)$ 换元，$\dd t'=(m_0c/F)\dd q'$：
\[
 x(t)=\int_0^t u(t')\dd t'
 =\frac{m_0c^2}{F}\int_0^{q(t)}
       \frac{q'\,\dd q'}{\sqrt{1+q'^2}}.
\]
再令 $z=1+q'^2$，$\dd z=2q'\dd q'$；积分限由 $0,q(t)$ 变成 $1,1+q(t)^2$，
故
\[
 \int_0^{q(t)}\frac{q'\dd q'}{\sqrt{1+q'^2}}
 =\frac12\int_1^{1+q(t)^2}z^{-1/2}\dd z
 =[\sqrt z]_1^{1+q(t)^2}
 =\sqrt{1+q(t)^2}-1.
\]
代回得到
\[
 x(t)=\frac{m_0c^2}{F}\left[\sqrt{1+\left(\frac{Ft}{m_0c}\right)^2}-1\right],
\]
故 $Fx(t)=m_0c^2(\sqrt{1+q^2}-1)=K(t)$。图\ref{fig:hr11s}两条曲线虽分别表示速度和动能，却由同一条粒子运动联系起来。若误用 $u=Ft/m_0$ 一直外推，既不能保持图中的速度上限，也不能与该功关系相容。
\end{solution}

\begin{question}
\begin{example}[康普顿散射的波长变化]\label{ex:hr10}
波长为 $\lambda_0$ 的光子与不变质量为 $m$ 的静止自由粒子碰撞。光子碰后相对入射方向偏转 $\theta$，见图\ref{fig:hr10q}。求出射光子的波长 $\lambda$。
\end{example}
\fig{fig:hr10q}{光子向右入射，碰后向右上方散射，粒子向右下方反冲。$\theta$ 是两光子传播方向之间的夹角。}{
\begin{tikzpicture}
\draw[light](-3,0)--(0,0);\fill(0,0)circle(2pt)node[below left]{$m$};\draw[axisline](0,0)--(3.7,0);
\draw[light](0,0)--(60:2.8)node[above]{$\lambda$};\draw[axis](0,0)--(-40.893:3)node[right]{反冲粒子};\node[above]at(-1.8,0){$\lambda_0$};\draw(.75,0)arc(0:60:.75);\node at(.95,.45){$\theta$};
\end{tikzpicture}}
\end{question}

\begin{solution}\textbf{根据动量三角形消去反冲方向}\par\nobreak
记入射、出射光子能量为 $\varepsilon_0=hc/\lambda_0$、$\varepsilon=hc/\lambda$，反冲粒子的能量、动量为 $E_e,\bm p_e$。图\ref{fig:hr10s}用首尾相接的矢量表示动量守恒：
\[
\bm p_e=\bm p_0-\bm p_\gamma,\qquad
p_e^2c^2=\varepsilon_0^2+\varepsilon^2-2\varepsilon_0\varepsilon\cos\theta.
\]


\par\medskip\Needspace{4\baselineskip}\noindent\textbf{把反冲总能量写成光子能量之差}\par\nobreak
能量守恒给出 $E_e=mc^2+\varepsilon_0-\varepsilon$。
\fig{fig:hr10s}{光子动量与反冲动量组成矢量三角形。图示长度取 $\varepsilon_0=mc^2,\theta=60^\circ$，两光子动量夹角为 $\theta$。}{
\begin{tikzpicture}
\coordinate(O)at(0,0);\coordinate(G)at(1.333333,2.309401);\coordinate(P)at(4,0);
\draw[axis](O)--(P)node[midway,below]{$\bm p_0$};\draw[axis](O)--(G)node[midway,above left]{$\bm p_\gamma$};\draw[axis](G)--(P)node[midway,above right]{$\bm p_e$};\draw(.7,0)arc(0:60:.7);\node at(.9,.45){$\theta$};\node[below left]at(O){$O$};
\end{tikzpicture}}



\par\medskip\Needspace{4\baselineskip}\noindent\textbf{联立反冲粒子的能量动量关系}\par\nobreak
将 $E_e$ 代入粒子的能量动量关系 $E_e^2-p_e^2c^2=m^2c^4$，有
\[
(mc^2+\varepsilon_0-\varepsilon)^2
-\left(\varepsilon_0^2+\varepsilon^2-2\varepsilon_0\varepsilon\cos\theta\right)
=m^2c^4.
\]


\par\medskip\Needspace{4\baselineskip}\noindent\textbf{整理出波长变化并检查两个极限}\par\nobreak
消去相同项，整理为
\[
mc^2(\varepsilon_0-\varepsilon)
=\varepsilon_0\varepsilon(1-\cos\theta).
\]
两边除以 $\varepsilon_0\varepsilon$，再换成波长，得到
\[
\boxed{\lambda-\lambda_0=\frac h{mc}(1-\cos\theta).}
\]
前向散射 $\theta=0$ 时波长不变；反向散射 $\theta=\pi$ 时增量最大，为 $2h/(mc)$。图\ref{fig:hr10q}中的反冲动量带走能量，所以出射光子一般波长增大；这里要求粒子起初静止且可视为自由粒子。



\par\medskip\Needspace{4\baselineskip}\noindent\textbf{计算反冲分量核对图中的实际角度}\par\nobreak
为了进一步核对图\ref{fig:hr10s}的矢量长度，可把结果写成出射能量：
\[
\varepsilon=\frac{\varepsilon_0}{1+\dfrac{\varepsilon_0}{mc^2}(1-\cos\theta)},\qquad
K_e=\varepsilon_0-\varepsilon.
\]
反冲动量的分量为 $p_{ex}=(\varepsilon_0-\varepsilon\cos\theta)/c$、$p_{ey}=-\varepsilon\sin\theta/c$，故向下反冲角 $\phi$ 满足
\[
\tan\phi=\frac{\varepsilon\sin\theta}{\varepsilon_0-\varepsilon\cos\theta}.
\]
图示参数给出 $\varepsilon=2\varepsilon_0/3$、$\phi=40.89^\circ$，与题设图的反冲方向和动量三角形一致。出射光子的 $\theta$ 与粒子反冲角 $\phi$ 是两个角，进入波长公式的是前者。
\end{solution}



\begin{question}
\begin{example}[寿命、飞行距离与同时性]\label{ex:r61}
某不稳定粒子的固有平均寿命为 \(\tau_0=2.00\,\mu\mathrm s\)，相对地面匀速 \(v=0.800c\)。取产生事件 \(A\) 为两系共同原点，求平均意义下的地面寿命与飞行距离。再在粒子静止系中用地面两个标记点的运动解释同一结果。取 \(c=3.00\times10^8\,\mathrm{m/s}\)。
\end{example}
\fig{fig:r61q}{\(B\) 表示经过一个平均寿命时的代表事件；平均寿命不是每个粒子的固定衰变时刻}{
\begin{tikzpicture}
\draw[axis](0,0)--(5,0)node[right]{$x$};
\draw[axis](0,0)--(0,5.7)node[above]{$ct$};
\draw[world](0,0)--(4,5);\draw[aux](4,0)--(4,5.3);
\fill(0,0)circle(1pt)node[below left]{$A$};
\fill(4,5)circle(1pt)node[right]{$B$};
\node[rotate=51,above]at(2,2.5){粒子世界线};
\node[below]at(4,0){地面终点 \(x=L\)};
\end{tikzpicture}}
\end{question}

\begin{solution}\textbf{由粒子固有平均寿命求地面寿命和距离}\par\nobreak
图\ref{fig:r61q}的 \(A,B\) 在粒子系同一空间点，所以其时间间隔为固有平均寿命。由 \(\gamma=1/\sqrt{1-0.8^2}=5/3\)，有
\[
 \boxed{\Delta t=\gamma\tau_0=\frac{10}{3}\,\mu\mathrm s},
 \qquad\boxed{L=v\Delta t=800\,\mathrm m}.
\]


\par\medskip\Needspace{4\baselineskip}\noindent\textbf{逐事件变换确认粒子系时间}\par\nobreak
将两个事件列在地面系：\(A:(0,0)\)，\(B:(L,\Delta t)\)。变换给出
\[
 \Delta x'=\gamma(L-v\Delta t)=0,\qquad
 \Delta t'=\gamma\left(\frac{10}{3}-\frac{0.8\times800}{300}\right)
 \mu\mathrm s=2.00\,\mu\mathrm s.
\]
\fig{fig:r61s}{\(A,D\) 在粒子系同时；\(A\) 与地面终点在 \(t=0\) 的事件并不在粒子系同时}{
\begin{tikzpicture}
\draw[axis](0,0)--(5,0)node[right]{$x$};
\draw[axis](0,0)--(0,5.7)node[above]{$ct$};
\draw[world](0,0)--(4,5);
\draw[world](4,0)--(4,5.4);
\draw[dashed](0,0)--(4.6,3.68)node[right]{$t'=0$};
\fill(4,3.2)circle(1pt)node[right]{$D$};
\fill(4,5)circle(1pt)node[right]{$B$};
\node[below left]at(0,0){$A$};\node[below]at(4,0){$L$};
\draw[aux](0,5)--(4,5)node[midway,above]{$ct=1000\,\mathrm m$};
\end{tikzpicture}}



\par\medskip\Needspace{4\baselineskip}\noindent\textbf{选粒子系同时事件测地面标记间距}\par\nobreak
图\ref{fig:r61s}中，\(D\) 是终点世界线与 \(t'=0\) 线的交点。地面原点与终点之间的固有距离为 \(L\)，在粒子系同时测得其间距
\[
 L'=\frac L\gamma=480\,\mathrm m.
\]


\par\medskip\Needspace{4\baselineskip}\noindent\textbf{用终点运动复核粒子系飞行时间}\par\nobreak
这段地面装置以 \(-0.800c\) 运动，终点从距粒子 \(480\,\mathrm m\) 处移到粒子旁。所用时间为
\[
 \frac{L'}{v}=\frac{480\,\mathrm m}{0.800c}=2.00\,\mu\mathrm s,
\]
与固有寿命一致。这里 \(L'\) 使用 \(A,D\) 两个同时事件；若直接变换地面 \(t=0\) 的两个端点事件，会得到另一对不同步事件，不能据此测粒子系长度。
\end{solution}

\begin{question}
\begin{example}[两个光子的系统质量]\label{ex:r62}
实验室中两光子的能量分别为 \(E_1=3.00\,\mathrm{MeV},E_2=1.00\,\mathrm{MeV}\)，传播夹角为 \(\theta\)。求系统不变质量随 \(\theta\) 的关系。对迎头相向的 \(\theta=\pi\)，求质心系相对实验室的速度及质心系中两光子的能量。
\end{example}
\fig{fig:r62q}{两光子的动量必须作矢量相加；单个光子仍有零静质量}{
\begin{tikzpicture}
\draw[axis,thick](0,0)--(3,0)node[right]{$\boldsymbol p_1$};
\draw[axis,thick](0,0)--(.5,.866025)node[above]{$\boldsymbol p_2$};
\draw(0.7,0)arc[start angle=0,end angle=60,radius=.7];
\node at(.9,.5){$\theta$};\node[below left]at(0,0){$O$};
\end{tikzpicture}}
\end{question}

\begin{solution}\textbf{将两光子动量作矢量相加}\par\nobreak
按图\ref{fig:r62q}，\(|\boldsymbol p_j|=E_j/c\)，所以
\[
 c^2|\boldsymbol p_1+\boldsymbol p_2|^2
 =E_1^2+E_2^2+2E_1E_2\cos\theta.
\]


\par\medskip\Needspace{4\baselineskip}\noindent\textbf{求任意夹角下的系统不变质量}\par\nobreak
代入总能量得
\[
 \boxed{M^2c^4=2E_1E_2(1-\cos\theta)},\qquad
 \boxed{Mc^2=2\sqrt{E_1E_2}\,|\sin(\theta/2)|}.
\]
\fig{fig:r62s}{迎头时实验室总动量向右；质心系中两光子等能、反向}{
\begin{tikzpicture}
\node[left]at(0,1){实验室};
\draw[axis,thick](.5,1)--(3.5,1)node[midway,above]{$E_1=3$};
\draw[axis,thick](5,1)--(4,1)node[midway,above]{$E_2=1$};
\node[left]at(0,-1){质心系};
\draw[axis,thick](.5,-1)--(2.232,-1)node[midway,above]{$E'_1=\sqrt3$};
\draw[axis,thick](5,-1)--(3.268,-1)node[midway,above]{$E'_2=\sqrt3$};
\node[below]at(2.7,-1.5){能量单位：\(\mathrm{MeV}\)};
\end{tikzpicture}}



\par\medskip\Needspace{4\baselineskip}\noindent\textbf{列出迎头情形的总能量和总动量}\par\nobreak
取图\ref{fig:r62s}向右为正，迎头时
\[
 E_{\rm tot}=4.00\,\mathrm{MeV},\quad
 cp_{\rm tot}=2.00\,\mathrm{MeV},\quad
 \boxed{Mc^2=2\sqrt3\,\mathrm{MeV}}.
\]


\par\medskip\Needspace{4\baselineskip}\noindent\textbf{求使系统总动量消失的参考系速度}\par\nobreak
要求质心系总动量为零，由 \(p'_{\rm tot}=\gamma_V(p_{\rm tot}-VE_{\rm tot}/c^2)=0\)，得
\[
 \boxed{V=\frac{c^2p_{\rm tot}}{E_{\rm tot}}=+0.500c}.
\]


\par\medskip\Needspace{4\baselineskip}\noindent\textbf{变换两光子并检查同向传播边界}\par\nobreak
再分别变换光子能量，\(\gamma_V=2/\sqrt3\)，
\[
 E'_1=\gamma_VE_1(1-0.5)=\sqrt3\,\mathrm{MeV},\qquad
 E'_2=\gamma_VE_2(1+0.5)=\sqrt3\,\mathrm{MeV}.
\]
两光子反向等能，总动量确为零。若 \(\theta=0\)，则 \(M=0\)，系统整体同向传播，不能用小于光速的惯性系把它变成总动量为零；这是上述“质心静止系”存在条件的边界。
\end{solution}

\section{本讲小结}

事件的位置和时刻必须由同一参考系的尺与同步钟定义。洛伦兹变换保持时空间隔不变，同时改变空间和时间坐标；收到光信号的时刻还包含传播延迟。类时关联的因果先后不因换系颠倒，类空事件的同时性和先后则可以依赖参考系。

时间膨胀比较同一只钟经历的两个事件，长度收缩要求在测量系同时记录两端。换系后，为测量长度通常须重选端点事件。速度变换把同一粒子相邻事件的空间差与时间差同时变换，因此不能直接相加高速速度，也不能把横向坐标不变误解为横向速度不变。

有质量粒子的总能量和动量满足 $E^2-p^2c^2=m_0^2c^4$。做功决定动能增量，相对论关系决定由动能反求的速度；恒力使动量线性增长，速度仍小于光速。碰撞、衰变与散射须在同一参考系内列出全部参与者的能量和动量，再由总量求系统不变质量。系统质量包含内部能量，通常不能逐个相加粒子的静质量。

光学实验中的条纹比较仍由往返传播时间与相位决定，介质中的光速须先注明介质静止系。局部加速效应与引力效应的等价有空间和时间范围限制；弱静态引力场中，时钟比较取决于引力势差。运动学效应与引力效应使用的条件不同，具体比较应把所选时钟、事件和测量量一并注明。

\section{练习}

\begin{question}
\begin{drill}[同一对闪光事件的同时性]\label{d:61}
飞船固有长度 \(300\,\mathrm m\)，沿地面 \(+x\) 方向以 \(0.600c\) 运动。地面系在同一时刻记录船尾闪光 \(A\) 和船头闪光 \(B\)。以 \(A\) 为两系原点，求 \(B\) 相对 \(A\) 的船系时间间隔，并判断先后。取 \(c=3.00\times10^8\,\mathrm{m/s}\)。
\end{drill}
\fig{fig:d61q}{练习：同一对闪光事件的同时性}{
\begin{tikzpicture}
\draw[axis](-.5,0)--(5,0)node[right]{$x$};
\draw[world](0,.4)rectangle(3.2,1.2);
\node[left]at(0,.8){$A$};\node[right]at(3.2,.8){$B$};
\draw[axis](1.1,1.6)--(2.6,1.6)node[midway,above]{$0.600c$};
\node[below]at(1.6,0){地面同时记录两端闪光};
\end{tikzpicture}}
\end{question}


\begin{question}
\begin{drill}[飞船中的横向运动]\label{d:62}
飞船相对地面沿 \(+x\) 方向以 \(V=0.600c\) 运动。在船系中，一粒子沿 \(+y'\) 方向以 \(0.800c\) 发射，且 \(u'_x=0\)。求地面系中粒子的两个速度分量与速率。
\end{drill}
\fig{fig:d62q}{练习：飞船中的横向运动}{
\begin{tikzpicture}
\draw[axis](0,0)--(3,0)node[right]{$x'$};
\draw[axis](0,0)--(0,2.7)node[above]{$y'$};
\draw[axis,thick](0,0)--(0,2)node[left]{$u'_y=0.800c$};
\draw[axis](4,0)--(6,0)node[midway,above]{$V=0.600c$};
\node[below]at(5,-.3){船相对地面};
\end{tikzpicture}}
\end{question}


\begin{question}
\begin{drill}[静止靶与运动粒子的粘合]\label{d:63}
两粒子静质量均为 \(m\)。实验室中一粒子以 \(0.600c\) 撞向另一静止粒子，碰后粘合且无能量逸出。求复合粒子的速度、不变质量，以及在复合粒子静止系中增加的内能。
\end{drill}
\fig{fig:d63q}{练习：静止靶与运动粒子的粘合}{
\begin{tikzpicture}
\draw[axisline](-1,0)--(6,0);
\fill(0,0)circle(2pt)node[below]{$m$};\fill(4,0)circle(2pt)node[below]{$m$};
\draw[axis](.4,.5)--(2.4,.5)node[midway,above]{$0.600c$};
\node[above]at(4,.3){静止};
\end{tikzpicture}}
\end{question}


\begin{question}
\begin{drill}[康普顿散射的能量与反冲方向]\label{d:64}
静止自由电子质量为 \(m\)。入射光子波长为 \(\lambda_C=h/(mc)\)，沿 \(+x\) 方向传播，散射后光子沿 \(+y\) 方向传播。求散射波长、电子动能和电子速度的两个分量。
\end{drill}
\fig{fig:d64q}{练习：康普顿散射的能量与反冲方向}{
\begin{tikzpicture}
\draw[ray](-2,0)--(0,0);\draw[ray](0,0)--(0,2);
\fill(0,0)circle(2pt)node[below right]{$e^-$};
\node[above]at(-1.4,0){$\lambda_C$};\node[right]at(0,1.4){散射光子};
\draw[axis](.4,-.5)--(2,-.5)node[right]{$+x$};
\end{tikzpicture}}
\end{question}


\begin{question}
\begin{drill}[飞船上观察百米赛跑]\label{hw:hr1}
飞船以 $V=0.8c$ 向北飞行。地面上一名运动员由南向北匀速跑完 $100\,\mathrm m$，地面记录的用时为 $10\,\mathrm s$。求飞船系中的跑道长度、运动员在这段运动中的路程及平均速度。取 $c=3.0\times10^8\,\mathrm{m/s}$，以北为正方向。
\end{drill}
\fig{fig:hrh1q}{地面系中的跑道与两种运动}{
\begin{tikzpicture}[x=1cm,y=1cm]
\draw[axis](-.4,0)--(8,0)node[right]{北 $+x$};
\draw[world](0,-.3)--(0,.3)node[above]{起点 $A$};\draw[world](5,-.3)--(5,.3)node[above]{终点 $B$};
\draw[dim](0,-.7)--(5,-.7)node[midway,fill=white]{$L_0=100\,\mathrm m$};
\fill(1.1,.15)circle(2pt);\draw[axis](1.1,.5)--(3.1,.5)node[midway,above]{$u=10\,\mathrm{m/s}$};
\draw[world](.4,1.5)rectangle(2,1.9);\draw[axis](2,1.7)--(4.4,1.7)node[midway,above]{$V=0.8c$};\node[left]at(.4,1.7){飞船};
\end{tikzpicture}}
\end{question}


\begin{question}
\begin{drill}[两种安装方式的光钟]\label{hw:hr2}
光钟由发光接收器 $A$ 与平面镜 $M$ 组成，二者在光钟静止系中的间距为 $l_0$，一次往返为一个滴答，固有周期 $\tau_0=2l_0/c$。分别讨论：
\begin{enumerate}[label=(\arabic*),leftmargin=2em]
\item 光钟固定在地面，$AM$ 垂直于车辆速度 $V$。求车上观察者测得的滴答周期。
\item 光钟固定在车上，$AM$ 平行于 $V$。求地面测得的 $AM$ 间距。
\item 对第 (2) 种安装方式，求地面测得的滴答周期。
\end{enumerate}
取 $c=3.0\times10^8\,\mathrm{m/s}$、$V=7.70\times10^7\,\mathrm{m/s}$、$l_0=2.46\times10^5\,\mathrm m$。
\end{drill}
\fig{fig:hrh2q}{分别固定在地面和车上的光钟}{
\begin{tikzpicture}[scale=.87]
\draw[world](0,0)--(0,2);\fill(0,0)circle(2pt)node[below]{$A$};\draw[world](-.7,2)--(.7,2)node[right]{$M$};\draw[dim](-1,0)--(-1,2)node[midway,left]{$l_0$};
\draw[world](-1.6,-.5)--(1.8,-.5);\node at(0,2.7){(1) 地面光钟};\draw[axis](.4,.8)--(2,.8)node[right]{$V$（车辆）};
\begin{scope}[shift={(5,0)}]
\draw[world](-.5,-.5)rectangle(3.5,1.5);\fill(0,.5)circle(2pt)node[below]{$A$};\draw[world](2.5,-.1)--(2.5,1.1)node[above]{$M$};
\draw[dim](0,1.9)--(2.5,1.9)node[midway,fill=white]{$l_0$（车上测量）};\draw[axis](3.7,.5)--(5,.5)node[right]{$V$};\node at(1.2,2.7){(2)、(3) 车载光钟};
\end{scope}
\end{tikzpicture}}
\end{question}


\begin{question}
\begin{drill}[恒星光行差]\label{hw:hr3}
航天飞机沿 $+x$ 方向飞行，速率为 $V$。在恒星静止系 $S$ 中，从飞机指向恒星的视线与 $+x$ 轴成 $\theta$ 角。求飞机系 $S'$ 中的视线角 $\theta'$。若前端半球形观察室只能接收 $\theta'<\pi/2$ 的恒星，说明 $V\to c$ 时可见范围怎样变化。取 $V=2.61\times10^8\,\mathrm{m/s}$、$c=3.0\times10^8\,\mathrm{m/s}$、$\theta=30.9^\circ$。
\end{drill}
\fig{fig:hrh3q}{视线指向恒星，光的传播方向与视线相反}{
\begin{tikzpicture}[scale=1]
\coordinate(O)at(0,0);\coordinate(H)at(4.8,2.872);\draw[axis](-.7,0)--(6,0)node[right]{$+x$};
\draw[aux](O)--(H);\draw[light](H)--(O);\fill(H)circle(2pt)node[above]{恒星};\fill(O)circle(2pt)node[below]{飞机};
\draw[axis](.3,-.6)--(2.5,-.6)node[right]{$V$};\draw(1,0)arc(0:30.9:1)node[midway,right]{$\theta$};\node at(3.1,2.5){入射光};
\end{tikzpicture}}
\end{question}


\begin{question}
\begin{drill}[高速粒子的完全非弹性碰撞]\label{hw:hr4}
两个相同质点的静质量均为 $m_0$，在实验室系中以等值反向的速率 $v$ 相撞，碰后粘成一个质点，过程中没有能量逸出。求复合质点的静质量。取 $v=2.74\times10^8\,\mathrm{m/s}$、$c=3.0\times10^8\,\mathrm{m/s}$、$m_0=4.21\times10^{-27}\,\mathrm{kg}$。
\end{drill}
\fig{fig:hrh4q}{两粒子的动量大小相等、方向相反}{
\begin{tikzpicture}
\draw[axis](-4,0)--(4,0)node[right]{$x$};\fill(-2,0)circle(3pt)node[below]{$m_0$};\fill(2,0)circle(3pt)node[below]{$m_0$};
\draw[axis](-2,.5)--(-.5,.5)node[midway,above]{$v$};\draw[axis](2,.5)--(.5,.5)node[midway,above]{$v$};\node[below]at(0,-.4){实验室系 $S$};
\end{tikzpicture}}
\end{question}


\begin{question}
\begin{drill}[两光子反推母粒子的静质量]\label{hw:hr5}
一个粒子衰变为两个光子。在同一惯性系中测得它们的能量分别为 $E_1,E_2$，传播方向夹角为 $\theta$。求母粒子的静质量 $m_0$。取 $E_1=182\,\mathrm{MeV}$、$E_2=162\,\mathrm{MeV}$、$\theta=37.6^\circ$。
\end{drill}
\fig{fig:hrh5q}{两光子的传播方向与夹角}{
\begin{tikzpicture}
\fill(0,0)circle(3pt)node[below left]{衰变点};\draw[light](0,0)--(4.8,0)node[right]{$E_1,\ \bm p_1$};\draw[light](0,0)--(37.6:{4.8*162/182})node[right]{$E_2,\ \bm p_2$};
\draw(1,0)arc(0:37.6:1)node[midway,right]{$\theta$};
\end{tikzpicture}}
\end{question}


\begin{question}
\begin{drill}[发射光子后的反冲]\label{hw:hr6}
静质量为 $M_0$ 的粒子起初静止，发射一个能量为 $Q=\lambda M_0c^2$ 的光子后，剩余粒子反冲。求反冲速度与剩余粒子的静质量。取 $\lambda=0.139$。
\end{drill}
\fig{fig:hrh6q}{初态静止，末态的光子和剩余粒子反向运动}{
\begin{tikzpicture}
\fill(-3.7,0)circle(4pt)node[below]{$M_0,\ u_0=0$};\draw[axis](-2.7,0)--(-1.4,0)node[midway,above]{发射};
\fill(1,0)circle(3pt)node[below]{$m$};\draw[axis](1,.5)--(-.6,.5)node[midway,above]{$u$};\draw[light](1,0)--(4,0)node[right]{$Q,\ p_\gamma=Q/c$};
\end{tikzpicture}}
\end{question}


\LectureAnswerSection

\begin{drillsolution}{d:61}\textbf{先在地面同时测量船长}\par\nobreak
图\ref{fig:d61q}两端同时测得的地面船长为
\[
 \gamma=1.25,\qquad L=300/1.25=240\,\mathrm m.
\]


\par\medskip\Needspace{4\baselineskip}\noindent\textbf{对同一对闪光变换时间差}\par\nobreak
固定事件顺序为 \(B-A\)，地面事件间隔为 \((\Delta x,\Delta t)=(240\,\mathrm m,0)\)，所以
\[
 \boxed{\Delta t'=-\gamma\frac{VL}{c^2}=-0.600\,\mu\mathrm s}.
\]
故船系中 \(\boxed{\text{船头闪光比船尾早 }0.600\,\mu\mathrm s}\)。


\par\medskip\Needspace{4\baselineskip}\noindent\textbf{逐个变换并核对另一条同时线}\par\nobreak
取事件 \(A\) 为两系共同原点，逐个变换事件可作独立核对：
\begin{center}\begin{tabular}{c|rr|rr}
事件&\(x/\mathrm m\)&\(t/\mu\mathrm s\)&\(x'/\mathrm m\)&\(t'/\mu\mathrm s\)\\\hline
船尾 \(A\)&0&0&0&0\\
船头 \(B\)&240&0&300&\(-0.600\)
\end{tabular}\end{center}
两系比较的是同一对闪光事件；船系的 \(300\,\mathrm m\) 也是两端静止位置之差，不能把时间差随之误写成零。
图\ref{fig:d61s}中，船头世界线与 $t'=0$ 相交于 $C$，闪光事件 $B$ 位于 $C$ 之前。$A,C$ 是另一对船系同时的端点事件；题目所问仍是 $A,B$，不能为恢复同时性而擅自换成 $C$。



\par\medskip\Needspace{4\baselineskip}\noindent\textbf{检查先后变化与因果关系}\par\nobreak
两事件类空分离，先后次序依赖参考系；这不构成从一端到另一端的超光速信号。
\fig{fig:d61s}{同一对闪光 $A,B$ 在地面同时；$C$ 才是船头在船系 $t'=0$ 上的事件}{\begin{tikzpicture}

\draw[axis](-.3,0)--(5.2,0)node[right]{$x/(100\,\mathrm m)$};
\draw[axis](0,-.5)--(0,4.1)node[above]{$ct/(100\,\mathrm m)$};
\draw[world](0,0)--(2.1,3.5)node[above] {船尾};
\draw[world](2.4,0)--(4.5,3.5)node[above]{船头};
\draw[dash dot,blue](0,0)--(4.7,2.82)node[right]{$t'=0$};
\draw[thick](0,0)--(2.4,0);
\fill(0,0)circle(1.5pt)node[below left]{$A$};
\fill(2.4,0)circle(1.5pt)node[below]{$B$};
\fill[blue](3.75,2.25)circle(1.5pt)node[above left]{$C$};
\node[below]at(1.2,-.65){地面同时：$t=0$};

\end{tikzpicture}}
\end{drillsolution}

\begin{drillsolution}{d:62}\textbf{取粒子世界线上相邻事件}\par\nobreak
图\ref{fig:d62q}的粒子在船系纯横向运动。为避免把速度分量当作普通矢量直接平移，选取粒子世界线上相邻两个事件，则 \(\dd x'=0\)、\(\dd y'=u'_y\dd t'\)。由逆洛伦兹变换，
\[
 \dd t=\gamma_V\left(\dd t'+\frac{V\dd x'}{c^2}\right)=\gamma_V\dd t',
 \quad \dd x=\gamma_VV\dd t',\quad \dd y=\dd y'.
\]


\par\medskip\Needspace{4\baselineskip}\noindent\textbf{由空间与时间微分共同求速度分量}\par\nobreak
因此 \(u_x=\dd x/\dd t=V\)，\(u_y=\dd y/\dd t=u'_y/\gamma_V\)。横向空间间隔不变，横向速度仍因时间间隔变换而改变。这里 \(\gamma_V=1.25\)、\(u'_x=0\)，故
\[
 \boxed{u_x=0.600c},\qquad
 \boxed{u_y=0.800c/1.25=0.640c}.
\]


\par\medskip\Needspace{4\baselineskip}\noindent\textbf{合成总速率并检查小于光速}\par\nobreak
图\ref{fig:d62s}中两分量组成地面系速度矢量，
\[
 \boxed{u=c\sqrt{0.600^2+0.640^2}=0.877c<c}.
\]
若把横向分量仍写成 \(0.800c\)，就遗漏了时间变换，算得速率恰为 \(c\)，与有静质量粒子的运动不符。
\fig{fig:d62s}{练习\ref{d:62}的解答图：地面系速度矢量由重新变换后的两个分量组成}{
\begin{tikzpicture}
\draw[axis](0,0)--(3,0)node[right]{$u_x$};
\draw[axis](0,0)--(0,3)node[above]{$u_y$};
\draw[axis,thick](0,0)--(1.8,1.92)node[above right]{$\boldsymbol u$};
\draw[aux](1.8,0)--(1.8,1.92)--(0,1.92);
\node[below]at(1.8,0){$0.600c$};\node[left]at(0,1.92){$0.640c$};
\end{tikzpicture}}
\end{drillsolution}

\begin{drillsolution}{d:63}\textbf{把静止靶的静能列入初态总能量}\par\nobreak
图\ref{fig:d63q}中，静止靶虽无动量，仍贡献静能 $mc^2$。运动粒子的 \(\gamma=1.25\)，所以
\[
 E_{\rm tot}=2.25mc^2,\qquad p_{\rm tot}=0.750mc.
\]


\par\medskip\Needspace{4\baselineskip}\noindent\textbf{由总能量和动量同时求速度与质量}\par\nobreak
封闭系统的能量、动量均守恒。实验室总动量不为零，复合粒子仍运动，须同时使用 $V=c^2p_{\rm tot}/E_{\rm tot}$ 和系统质量不变量：
\[
 \boxed{V=\frac{c^2p_{\rm tot}}{E_{\rm tot}}=\frac c3},\qquad
 \boxed{M=m\sqrt{2.25^2-0.750^2}=\frac3{\sqrt2}m}.
\]
图\ref{fig:d63s}显示碰后实验室动量仍向右，不能把实验室总能量直接写成 \(Mc^2\)。


\par\medskip\Needspace{4\baselineskip}\noindent\textbf{在复合物静止系计算增加的内能}\par\nobreak
在复合粒子静止系中增加的内能为
\[
 \boxed{\Delta E_{\rm int}=(M-2m)c^2=(3/\sqrt2-2)mc^2\approx0.121mc^2}.
\]
这不是实验室初始全部动能 \(0.250mc^2\)，其余部分对应碰后整体运动。
\fig{fig:d63s}{练习\ref{d:63}的解答图：实验室仍有整体动量；静止系总能量才等于 \(Mc^2\)}{
\begin{tikzpicture}
\node at(0,1){实验室碰后};
\fill(2,1)circle(2pt)node[below]{$M$};\draw[axis](2.3,1)--(4.3,1)node[above]{$V=c/3$};
\node at(0,-1){复合粒子静止系};
\fill(2,-1)circle(2pt)node[below]{$M,\ \boldsymbol p'=0$};
\node[right]at(3,-1){$E'=Mc^2>2mc^2$};
\end{tikzpicture}}
\end{drillsolution}

\begin{drillsolution}{d:64}\textbf{用九十度散射求光子波长}\par\nobreak
图\ref{fig:d64q}中散射角 \(90^\circ\)，所以
\[
 \boxed{\lambda=\lambda_C+\lambda_C(1-\cos90^\circ)=2\lambda_C}.
\]


\par\medskip\Needspace{4\baselineskip}\noindent\textbf{由能量守恒求电子动能和总能量}\par\nobreak
初、末光子能量分别为 \(mc^2,mc^2/2\)。由能量守恒，
\(\boxed{K_e=mc^2/2}\)，电子总能量 \(E_e=3mc^2/2\)。


\par\medskip\Needspace{4\baselineskip}\noindent\textbf{按矢量差求电子动量分量}\par\nobreak
初态电子静止，图\ref{fig:d64s}中电子末动量为入射光子动量减去上行散射光子动量，所以
\[
 p_{ex}=mc,\qquad p_{ey}=-mc/2.
\]


\par\medskip\Needspace{4\baselineskip}\noindent\textbf{以总能量和动量共同反求电子速度}\par\nobreak
利用 \(\boldsymbol v=c^2\boldsymbol p/E_e\)，得
\[
 \boxed{v_x=2c/3,\qquad v_y=-c/3},\qquad
 |\boldsymbol v|=\sqrt5c/3<c.
\]
电子向右下方反冲，以补偿散射光子向上的动量。这里不能用 \(p=mv\) 反求速度。
\fig{fig:d64s}{练习\ref{d:64}的解答图：上行光子动量与下行电子动量的纵向分量抵消}{
\begin{tikzpicture}
\draw[axis](0,0)--(2,0)node[right]{$\boldsymbol p_{\gamma0}=mc\,\boldsymbol e_x$};
\draw[axis,thick](0,0)--(0,1)node[above]{$\boldsymbol p_\gamma$};
\draw[axis,thick](0,0)--(2,-1)node[below right]{$\boldsymbol p_e$};
\draw[aux](0,1)--(2,0)--(2,-1);
\node[left]at(0,.5){$mc/2$};\node[right]at(2,-.5){$-mc/2$};
\end{tikzpicture}}
\end{drillsolution}

\begin{drillsolution}{hw:hr1}\textbf{在飞船同时面上测量跑道长度}\par\nobreak
图\ref{fig:hrh1q}中的跑道静止于地面，故 $L_0$ 是固有长度。飞船必须在自身同一时刻测量跑道两端，所得长度为
\[
L'=\frac{L_0}{\gamma}=60\,\mathrm m,\qquad \gamma=\frac53.
\]


\par\medskip\Needspace{4\baselineskip}\noindent\textbf{为赛跑另选起跑和到达两个事件}\par\nobreak
运动员的起跑、到达是另外两个事件，记作 $E_A,E_B$。图\ref{fig:hrh1s}把它们与跑道长度的测量分开表示。
\fig{fig:hrh1s}{赛跑事件的变换；两侧坐标采用各自注明的尺度}{
\begin{tikzpicture}
\begin{scope}[x=.34cm,y=.2cm]
\draw[axis](-.5,0)--(10,0)node[right]{$x/(100\,\mathrm m)$};\draw[axis](0,-.4)--(0,13)node[above]{$t/\mathrm s$};
\draw[world](0,0)--(1,10);\fill(0,0)circle(1.7pt)node[below left]{$E_A$};\fill(1,10)circle(1.7pt)node[right]{$E_B(1,10)$};\node at(6,4){地面系 $S$};
\end{scope}
\begin{scope}[shift={(8,0)},x=.65cm,y=.14cm]
\draw[axis](-5,0)--(1,0)node[right]{$x'/(10^9\,\mathrm m)$};\draw[axis](0,-.5)--(0,21)node[above]{$t'/\mathrm s$};
\draw[world](0,0)--(-4,16.667);\fill(0,0)circle(1.7pt)node[below left]{$E'_A$};\fill(-4,16.667)circle(1.7pt)node[above]{$E'_B$};\node at(-3.3,6){飞船系 $S'$};
\end{scope}
\end{tikzpicture}}



\par\medskip\Needspace{4\baselineskip}\noindent\textbf{变换赛跑事件的位移和时间}\par\nobreak
以 $E_A$ 为共同零点，图中的 $E_B$ 满足 $\Delta x=100\,\mathrm m$、$\Delta t=10\,\mathrm s$。洛伦兹变换给出
\begin{align*}
\Delta x'&=\gamma(\Delta x-V\Delta t)
 =-3.999999833\times10^9\,\mathrm m,\\
\Delta t'&=\gamma\left(\Delta t-\frac{V\Delta x}{c^2}\right)
 =16.66666622\,\mathrm s.
\end{align*}
图\ref{fig:hrh1s}中 $E'_B$ 在 $E'_A$ 左侧，说明运动员在飞船系向南运动。



\par\medskip\Needspace{4\baselineskip}\noindent\textbf{求路程与带方向的平均速度}\par\nobreak
他做匀速直线运动，路程为位移的绝对值：
\[
\boxed{s'=|\Delta x'|\simeq4.00\times10^9\,\mathrm m}.
\]
平均速度为
\[
\boxed{u'=\frac{\Delta x'}{\Delta t'}
 =\frac{u-V}{1-uV/c^2}
 \simeq-2.399999964\times10^8\,\mathrm{m/s}}.
\]
跑道在飞船系中也不断向南移动，因此运动员走过的路程不能用收缩后的跑道长度 $60\,\mathrm m$ 代替。路程取正值，向南的方向由速度的负号表示。


\par\medskip\Needspace{4\baselineskip}\noindent\textbf{由运动终点线核对赛跑路程}\par\nobreak
还可在飞船系内直接复核。取起跑事件 $E'_A$ 为零点，飞船系 $t'=0$ 时终点线在 $x'=L'=60\,\mathrm m$，随后以速度 $-V$ 向南移动。

运动员到达终点线时，故有
\[
\Delta x'=L'-V\Delta t'.
\]
代入 $\Delta t'=16.66666622\,\mathrm s$，与事件变换所得位移一致。这说明图\ref{fig:hrh1s}的长位移主要来自跑道随地面向后运动，并非把一段百米跑道测成了数十亿米。

\end{drillsolution}

\begin{drillsolution}{hw:hr2}\textbf{从固有镜距和车速计算共同参数}\par\nobreak
先由图\ref{fig:hrh2q}中的镜距 $l_0$和车速 $V$算参数：
\[
\beta=0.256667,\quad \gamma=1.034661,\quad
\tau_0=1.640\times10^{-3}\,\mathrm s.
\]


\par\medskip\Needspace{4\baselineskip}\noindent\textbf{用横向光路求车上测得的地面钟周期}\par\nobreak
第 (1) 问中，光钟随地面相对于车辆向后运动。图\ref{fig:hrh2s}左侧的 $A_0,M_1,A_2$ 是车上观察者所见的发射、反射和接收位置。光钟横向尺寸不变，每半程均有
\[
\left(\frac{c\tau}{2}\right)^2=l_0^2+
\left(\frac{V\tau}{2}\right)^2,
\quad\boxed{\tau=\gamma\tau_0=1.69684\times10^{-3}\,\mathrm s}.
\]
\fig{fig:hrh2s}{左：车上所见地面光钟的路径；右：地面所见车载纵向光钟的三个事件}{
\begin{tikzpicture}[scale=.9]
\pgfmathsetmacro{\halfshift}{(7.7/30)/sqrt(1-(7.7/30)^2)*2.4}
\draw[light](0,0)--(-\halfshift,2.4)--({-2*\halfshift},0);\draw[axisline]({-2*\halfshift},0)--(0,0);\draw[dim](.5,0)--(.5,2.4)node[midway,right]{$l_0$};
\fill(0,0)circle(1.6pt)node[below]{$A_0$};\fill(-\halfshift,2.4)circle(1.6pt)node[above]{$M_1$};\fill({-2*\halfshift},0)circle(1.6pt)node[below]{$A_2$};
\draw[dim]({-2*\halfshift},-.7)--(0,-.7)node[midway,fill=white]{$V\tau$};
\begin{scope}[shift={(4.4,-.3)},scale=.85]
\pgfmathsetmacro{\be}{7.7/30}\pgfmathsetmacro{\ell}{2*sqrt(1-\be*\be)}\pgfmathsetmacro{\tm}{\ell/(1-\be)}\pgfmathsetmacro{\tt}{2*\ell/(1-\be*\be)}
\draw[axis](-.2,0)--(4,0)node[right]{$x$};\draw[axis](0,-.2)--(0,5.2)node[above]{$ct$};
\draw[world](0,0)--({\be*4.9},4.9)node[above right]{$A$};\draw[world](\ell,0)--({\ell+\be*4.9},4.9)node[above]{$M$};
\draw[light](0,0)--(\tm,\tm)--({\be*\tt},\tt);
\fill(0,0)circle(1.6pt)node[below left]{$E_0$};\fill(\tm,\tm)circle(1.6pt)node[right]{$E_1$};\fill({\be*\tt},\tt)circle(1.6pt)node[left]{$E_2$};
\draw[dim](0,-.55)--(\ell,-.55)node[midway,fill=white]{$L=l_0/\gamma$};
\end{scope}
\end{tikzpicture}}



\par\medskip\Needspace{4\baselineskip}\noindent\textbf{同时测量车载纵向光钟的镜距}\par\nobreak
第 (2) 问的光钟固定于车上，地面同时测量图\ref{fig:hrh2s}右侧两条世界线之间的水平距离，得
\[
\boxed{L=\frac{l_0}{\gamma}=2.37759\times10^5\,\mathrm m}.
\]


\par\medskip\Needspace{4\baselineskip}\noindent\textbf{分别求追赶和迎面传播的两半程}\par\nobreak
第 (3) 问沿右图的折线 $E_0E_1E_2$ 计算。

去程追赶前方镜面，回程迎向接收器，所以
\[
t_+=\frac{L}{c-V},\qquad t_-=\frac{L}{c+V}.
\]
两段时间不等，但总周期仍为
\[
\boxed{\tau=t_++t_-=\frac{2Lc}{c^2-V^2}
 =\gamma\frac{2l_0}{c}=1.69684\times10^{-3}\,\mathrm s}.
\]


\par\medskip\Needspace{4\baselineskip}\noindent\textbf{检查完整滴答与半程的事件区别}\par\nobreak
第 (1) 问的相邻滴答发生在同一地面接收器处，第 (3) 问则发生在同一车载接收器处；各自的固有周期均为 $\tau_0$。
图\ref{fig:hrh2s}右侧还可读出两个不相等的半程：
\[
t_+=1.06618\times10^{-3}\,\mathrm s,\qquad
t_-=0.630660\times10^{-3}\,\mathrm s.
\]
它们的比值为 $(c+V)/(c-V)$，相加才是实验室周期。

不能把“一个滴答整体变慢”进一步理解成纵向光钟的每个半程都等于 $\gamma l_0/c$：反射事件 $E_1$ 在另一面镜上，时间变换中还含有位置项。横向光钟的两半程对称，纵向光钟则不对称，二者的完整周期才满足同一时间膨胀关系。

\end{drillsolution}

\begin{drillsolution}{hw:hr3}\textbf{由视线方向反向确定入射光速度}\par\nobreak
图\ref{fig:hrh3q}中的光向左下方传播，故在 $S$ 中
\[
u_x=-c\cos\theta,\qquad u_y=-c\sin\theta.
\]


\par\medskip\Needspace{4\baselineskip}\noindent\textbf{变换光速分量后重新定义视线角}\par\nobreak
速度变换给出
\[
u'_x=-c\frac{\cos\theta+\beta}{1+\beta\cos\theta},\qquad
u'_y=-c\frac{\sin\theta}{\gamma(1+\beta\cos\theta)}.
\]
在飞机系中，视线仍取光传播的反方向。图\ref{fig:hrh3s}左侧据此画出 $\theta'$，所以
\[
\boxed{\cos\theta'=\frac{\cos\theta+\beta}{1+\beta\cos\theta}},\qquad
\sin\theta'=\frac{\sin\theta}{\gamma(1+\beta\cos\theta)}.
\]
取 $\beta=0.87$，得 $\boxed{\theta'=8.34^\circ}$。
\fig{fig:hrh3s}{左：前方恒星的视线偏向机头；右：恒星系中的可见角范围}{
\begin{tikzpicture}[scale=.9]
\draw[axis](0,0)--(4.5,0)node[right]{机头};\draw[aux](0,0)--(3.8,2.274)node[above]{$\theta=30.9^\circ$};
\draw[world](0,0)--(4.2,.615)node[above left]{$\theta'=8.34^\circ$};\draw(1,0)arc(0:30.9:1);\draw(1.6,0)arc(0:8.34:1.6);
\fill(0,0)circle(2pt);
\begin{scope}[shift={(7.8,0)}]
\fill[gray!12](0,0)--(-150.458:1.8)arc(-150.458:150.458:1.8)--cycle;
\draw[world](0,0)circle(1.8);\draw[axis](-2.1,0)--(2.3,0)node[right]{$+x$};\draw[aux](0,0)--(150.458:1.8);\draw[aux](0,0)--(-150.458:1.8);
\node[align=center]at(.55,-.75){可见方向};\node[align=center]at(-2.1,-1.5){后方盲区};\draw[aux](-1.6,-.5)--(-2.1,-1.2);\draw(.7,0)arc(0:150.458:.7);\node at(.1,1.15){$\theta_{\max}$};
\end{scope}
\end{tikzpicture}}



\par\medskip\Needspace{4\baselineskip}\noindent\textbf{由前半球条件求原星空可见范围}\par\nobreak
半球观察室接收 $\cos\theta'>0$ 的光。分母 $1+\beta\cos\theta$ 恒正，故恒星系中的可见条件为
\[
\cos\theta>-\beta,\qquad
0\le\theta<\theta_{\max}=\arccos(-\beta).
\]
本题 $\theta_{\max}=150.46^\circ$，右图的后方盲区已小于半球。当 $\beta\to1$ 时，$\theta_{\max}\to\pi$，除严格后方的方向外，几乎整个星空都可进入前端观察室。

图中画的是过飞行轴的截面；绕轴旋转即为空间角范围。


\par\medskip\Needspace{4\baselineskip}\noindent\textbf{以半角关系解释视线聚向机头}\par\nobreak
角度的压缩还可写成半角形式。由上面的正弦、余弦关系，
\[
\tan\frac{\theta'}2=\frac{\sin\theta'}{1+\cos\theta'}
=\sqrt{\frac{1-\beta}{1+\beta}}\tan\frac\theta2.
\]
对任一固定的 $\theta<\pi$，当 $\beta\to1$ 时均有 $\theta'\to0$，即星光视线聚向机头。

\par\medskip\Needspace{4\baselineskip}\noindent\textbf{计算方向立体角并明确所数范围}\par\nobreak
若按恒星系的方向立体角计算，轴对称可见区域的立体角微元为
$\dd\Omega=\sin\theta\,\dd\theta\,\dd\phi$。先积方位角，再积极角，
\begin{align*}
 \Omega_{\text{可见}}
 &=\int_0^{2\pi}\dd\phi\int_0^{\theta_{\max}}\sin\theta\,\dd\theta\\
 &=2\pi[-\cos\theta]_0^{\theta_{\max}}
 =2\pi(1-\cos\theta_{\max})=2\pi(1+\beta).
\end{align*}
占整天方向的比例为 $(1+\beta)/2$，本题为 $93.5\%$。

这个比例描述可见的方向范围，未对恒星的亮度或实际数量分布作假设。

\end{drillsolution}

\begin{drillsolution}{hw:hr4}\textbf{由对称初态动量判断碰后静止}\par\nobreak
图\ref{fig:hrh4q}两侧动量相消，系统总动量为零。碰后只有一个质点，因而复合质点在实验室系静止。它的总能量就是静能 $Mc^2$。
\fig{fig:hrh4s}{碰前等值反向动量相消；碰后只有一个复合粒子，静止于同一实验室系}{\begin{tikzpicture}

\fill(-2,1)circle(2.5pt);\fill(2,1)circle(2.5pt);
\draw[axis,thick](-2,1.5)--(-.4,1.5)node[midway,above]{$\bm p_1$};
\draw[axis,thick](2,1.5)--(.4,1.5)node[midway,above]{$\bm p_2=-\bm p_1$};
\node[below]at(-2,1){$E_1=\gamma m_0c^2$};\node[below]at(2,1){$E_2=\gamma m_0c^2$};
\draw[axis](0,.4)--(0,-.5);
\fill(0,-1)circle(4pt)node[right]{$M,\ u=0$};
\node[below]at(0,-1.5){$Mc^2=E_1+E_2$};

\end{tikzpicture}}


\par\medskip\Needspace{4\baselineskip}\noindent\textbf{把两份相对论总能量转为复合静能}\par\nobreak
对图\ref{fig:hrh4s}写能量守恒式：
\[
Mc^2=2\gamma m_0c^2,\qquad
\boxed{M=\frac{2m_0}{\sqrt{1-v^2/c^2}}
 =2.0677\times10^{-26}\,\mathrm{kg}}.
\]
数值上 $\gamma=2.45572$，所以 $M=4.91144m_0$。增加的静能 $(M-2m_0)c^2$ 等于碰前的总动能，并以复合系统的内能等形式保留。


\par\medskip\Needspace{4\baselineskip}\noindent\textbf{以动能转内能核对质量增量}\par\nobreak
把图\ref{fig:hrh4s}上方每个粒子的能量拆开，有 $\gamma m_0c^2=m_0c^2+K_1$，因此也可写成
\[
M=2m_0+\frac{2K_1}{c^2}.
\]
粘合后宏观动能为零，复合系统的内能相应增加，总能量守恒。

这里 $v$ 是每个粒子在实验室系的速率，求 $\gamma$ 时应使用它；若换成两粒子的相对速率，就改变了能量所属的参考系。

\end{drillsolution}

\begin{drillsolution}{hw:hr5}\textbf{由两光子矢量和确定母粒子动量}\par\nobreak
母粒子在所选惯性系中未必静止。图\ref{fig:hrh5s}把图\ref{fig:hrh5q}中的两个光子动量首尾相接，合动量 $\bm p=\bm p_1+\bm p_2$ 就是母粒子衰变前的动量。
\fig{fig:hrh5s}{动量三角形；平行虚线保留两光子的夹角}{
\begin{tikzpicture}[scale=.9]
\coordinate(O)at(0,0);\coordinate(P)at(3.64,0);\coordinate(Q)at(6.207,1.977);
\draw[axis](O)--(P)node[midway,below]{$\bm p_1$};\draw[axis](P)--(Q)node[midway,above left]{$\bm p_2$};\draw[axis](O)--(Q)node[midway,above]{$\bm p$};
\draw[aux](P)--(5.2,0);\draw(4.44,0)arc(0:37.6:.8)node[midway,right]{$\theta$};
\node[below left]at(O){$O$};\node[below]at(P){$P$};\node[right]at(Q){$Q$};
\end{tikzpicture}}



\par\medskip\Needspace{4\baselineskip}\noindent\textbf{联立能量和动量求母粒子质量}\par\nobreak
由图中的矢量和及光子关系 $p_i=E_i/c$，
\[
p^2=\frac{E_1^2+E_2^2+2E_1E_2\cos\theta}{c^2}.
\]
总能量 $E=E_1+E_2$。代入母粒子的能量动量关系，得
\begin{align*}
m_0^2c^4&=(E_1+E_2)^2-p^2c^2\\
 &=2E_1E_2(1-\cos\theta).
\end{align*}
因此
\[
\boxed{m_0=\frac{\sqrt{2E_1E_2(1-\cos\theta)}}{c^2}
 =110.67\,\mathrm{MeV}/c^2}.
\]


\par\medskip\Needspace{4\baselineskip}\noindent\textbf{检查能否把母粒子误当静止}\par\nobreak
若母粒子静止，两个光子应等能且反向。本题的能量和夹角不满足这一条件，故不能令总动量为零。



\par\medskip\Needspace{4\baselineskip}\noindent\textbf{反求运动状态验证质量不变量}\par\nobreak
由图\ref{fig:hrh5s}还可检查母粒子的运动状态。它的总能量为 $344\,\mathrm{MeV}$，故
\[
\frac{u}{c}=\frac{pc}{E}
=\sqrt{1-\left(\frac{m_0c^2}{E_1+E_2}\right)^2}
\simeq0.94683.
\]
若选到母粒子静止系，总动量才为零，两光子在那里反向且各有能量 $m_0c^2/2$。当前系的两份能量可以不同，但由它们和夹角求出的 $m_0$ 保持不变；这正是不变量方法能避免另外求母粒子速度的原因。

\end{drillsolution}

\begin{drillsolution}{hw:hr6}\textbf{按动量反向和能量分配确定末态总量}\par\nobreak
取光子方向为 $+x$。图\ref{fig:hrh6q}中的初态总动量为零，故剩余粒子的动量为 $p_x=-Q/c$。能量守恒给出
\[
E=M_0c^2-Q=(1-\lambda)M_0c^2.
\]
图\ref{fig:hrh6s}分别画出能量分配和动量平衡。剩余粒子的 $E$ 包含静能与反冲动能，不能直接以 $E/c^2$ 作为静质量。
\fig{fig:hrh6s}{能量分配与等值反向的动量}{
\begin{tikzpicture}[scale=.95]
\draw[world](0,1)rectangle(7,1.7);\draw[world](6.027,1)--(6.027,1.7);\node at(3,1.35){$(1-\lambda)M_0c^2$};\node at(6.52,1.35){$Q$};\node[left]at(0,1.35){能量};
\draw[axis](3.5,0)--(1,0)node[midway,below]{$\bm p$};\draw[axis](3.5,0)--(6,0)node[midway,below]{$\bm p_\gamma$};\node[above]at(3.5,0){$|\bm p|=|\bm p_\gamma|$};
\end{tikzpicture}}



\par\medskip\Needspace{4\baselineskip}\noindent\textbf{由剩余粒子的总能量求反冲速率}\par\nobreak
用 $\bm u=c^2\bm p/E$，反冲速率为
\[
\boxed{\frac{u}{c}=\frac{\lambda}{1-\lambda}=0.16144},
\]
方向与图中光子相反。

\par\medskip\Needspace{4\baselineskip}\noindent\textbf{求剩余静质量并检查允许参数}\par\nobreak
再由能量动量关系，
\begin{align*}
m^2c^4&=E^2-p^2c^2\\
 &=[(1-\lambda)^2-\lambda^2]M_0^2c^4,
\end{align*}
得
\[
\boxed{\frac{m}{M_0}=\sqrt{1-2\lambda}=0.849706}.
\]
若剩余粒子具有正静质量，需 $0<\lambda<1/2$。

本题满足该条件，且所得反冲速率小于光速。


\par\medskip\Needspace{4\baselineskip}\noindent\textbf{核对静能减少如何供给光子和反冲动能}\par\nobreak
静质量的减少也可逐项核对。剩余粒子的反冲动能为
\[
K=E-mc^2=\bigl[1-\lambda-\sqrt{1-2\lambda}\bigr]M_0c^2
\simeq0.0112942M_0c^2.
\]
因此 $(M_0-m)c^2=Q+K$：初态减少的静能既供给光子，也供给反冲动能。图\ref{fig:hrh6s}左侧大能量块表示剩余粒子的总能量，其中还应包含这份 $K$，不能把它全部当作静能。

\end{drillsolution}

\end{document}
